% Generated mechanically from frozen input canon/ko-kr/integration-20260908-r10-p03-register/inputs/p03/ko/hypercovering.tex.
% Do not edit this generated file; edit the bound locale source instead.

% OK, start here.
%

\setcounter{chapter}{24}
\chapter{초덮개}



\phantomsection
\label{hypercovering-section-phantom}


\section{서론}
\label{hypercovering-section-introduction}

\noindent
$\mathcal{C}$를 사이트라 하자. Sites, 정의 \ref{sites-definition-site}를 보라.
$X$를 $\mathcal{C}$의 대상이라 하자.
$\mathcal{C}$ 위의 아벨 층 $\mathcal{F}$가 주어졌을 때,
그 코호몰로지 군
$$
H^i(X, \mathcal{F}).
$$
을 계산하고자 한다. 일반적인 정의(Cohomology on Sites, 절
\ref{sites-cohomology-section-cohomology-sheaves})에 따르면,
단사 분해
$
0 \to \mathcal{F} \to \mathcal{I}^0 \to \mathcal{I}^1 \to \ldots
$
를 택하고
$$
H^i(X, \mathcal{F})
=
H^i(
\Gamma(X, \mathcal{I}^0) \to
\Gamma(X, \mathcal{I}^1) \to
\Gamma(X, \mathcal{I}^2)\to \ldots)
$$
로 놓아서 이 코호몰로지 군을 계산한다. 이 장의 목표는
$\mathcal{C}$가 섬유곱을 가질 때에는 단사 분해를 택하지 않고도
이 코호몰로지 군들을 계산할 수 있음을 보이는 것이다.
이를 위해 초덮개를 사용한다.

\medskip\noindent
사이트에서 초덮개는 덮개의 일반화이다. \cite[Expos\'e V, Sec. 7]{SGA4}을
보라. 대상 $X$의 초덮개 $K$가 주어지면, $X$ 위의 아벨 층
$\mathcal{F}$의 코호몰로지를 $K$의 성분 $K_n$ 위에서 그 층이 갖는
코호몰로지로 나타내는 체흐-코호몰로지 스펙트럼 열이 있다.
항상 충분히 많은 초덮개가 존재함이 드러난다. 따라서 모든 초덮개에
걸쳐 여극한을 취하면 스펙트럼 열이 퇴화하고, $X$ 위에서
$\mathcal{F}$의 코호몰로지는 체흐 코호몰로지 군들의 여극한으로 계산된다.

\medskip\noindent
더 일반적으로는 코호몰로지 강하를 갖는 단체적 확대를 고려할 수 있다.
\cite[Expos\'e Vbis]{SGA4}을 보라. 코호몰로지 강하를 다룬 좋은 원고로
Brian Conrad의 글이 있다. 다음을 보라.
\url{https://math.stanford.edu/~conrad/papers/hypercover.pdf}.
단체적 공간을 다루는 장에서 이 문제들로 돌아와, 예를 들어
``국소 콤팩트'' 위상공간의 고유 초덮개가 코호몰로지 강하를 이룸을
보일 것이다(Simplicial Spaces, 절
\ref{spaces-simplicial-section-proper-hypercovering}).
증명 전략은 이 명제를 이 장에서 구성하는 체흐-코호몰로지
스펙트럼 열로 귀착시키는 것이다.






























\section{반표현 가능 대상}
\label{hypercovering-section-semi-representable}

\noindent
먼저 다음 정의를 둔다.
문자 ``SR''은 Semi-Representable의 머리글자이다.

\begin{definition}
\label{hypercovering-definition-SR}
$\mathcal{C}$를 범주라 하자. 다음과 같이 정의되는
{\it 반표현 가능 대상}의 범주를 $\text{SR}(\mathcal{C})$로 나타낸다.
\begin{enumerate}
\item 대상은 대상들의 족 $\{U_i\}_{i \in I}$이고,
\item 사상 $\{U_i\}_{i \in I} \to \{V_j\}_{j \in J}$는 사상
$\alpha : I \to J$와 각 $i \in I$에 대한 $\mathcal{C}$의 사상
$f_i : U_i \to V_{\alpha(i)}$로 주어진다.
\end{enumerate}
$X \in \Ob(\mathcal{C})$를 $\mathcal{C}$의 대상이라 하자.
{\it $X$ 위의 반표현 가능 대상}의 범주는
$\text{SR}(\mathcal{C}, X) = \text{SR}(\mathcal{C}/X)$.
\end{definition}

\noindent
이 정의는 본질적으로
\cite[Expos\'e V, Subsection 7.3.0]{SGA4}의 정의와 동치이다.
이는 ``큰'' 범주임에 유의하자. 나중에 $X$의 초덮개에 필요한
첨자 집합 $I$의 크기를 ``제한''한다. 그러면
$\text{SR}(\mathcal{C}, X)$가 범주가 되도록 다시 정의할 수 있다.
$\text{SR}(\mathcal{C}, X)$의 대상과 사상을 명시하면 다음과 같다.
\begin{enumerate}
\item 대상은 사상들의 족 $\{U_i \to X\}_{i \in I}$이고,
\item 사상 $\{U_i \to X\}_{i \in I} \to
\{V_j \to X\}_{j \in J}$는 사상 $\alpha : I \to J$와
각 $i \in I$에 대한 $X$ 위의 사상
$f_i : U_i \to V_{\alpha(i)}$로 주어진다.
\end{enumerate}
망각 함자
$\text{SR}(\mathcal{C}, X) \to \text{SR}(\mathcal{C})$가 있다.

\begin{definition}
\label{hypercovering-definition-SR-F}
$\mathcal{C}$를 범주라 하자.
{\it 반표현 가능 대상에 준층을 대응시키는 함자}를 $F$로 나타낸다.
식으로 쓰면
\begin{eqnarray*}
F : \text{SR}(\mathcal{C}) & \longrightarrow & \textit{PSh}(\mathcal{C}) \\
\{U_i\}_{i \in I} & \longmapsto & \amalg_{i\in I} h_{U_i}
\end{eqnarray*}
이며, 여기서 $h_U$는 대상 $U$에 대응하는 표현 가능 준층을 나타낸다.
\end{definition}

\noindent
사상 $U \to X$가 주어지면 표현 가능 준층의 사상 $h_U \to h_X$를 얻는다.
따라서 흔히 $\text{SR}(\mathcal{C}, X)$ 위의 $F$를 $h_X$ 위의
집합 값 준층의 범주, 즉 $\textit{PSh}(\mathcal{C})/h_X$로 가는
함자로 생각한다. 다음은 그 도표이다.
$$
\xymatrix{
\text{SR}(\mathcal{C}, X) \ar[r]_F \ar[d] &
\textit{PSh}(\mathcal{C})/h_X \ar[d] \\
\text{SR}(\mathcal{C}) \ar[r]^F &
\textit{PSh}(\mathcal{C})
}
$$
이제 반표현 가능 대상의 범주에서 극한의 존재를 논한다.

\begin{lemma}
\label{hypercovering-lemma-coprod-prod-SR}
$\mathcal{C}$를 범주라 하자.
\begin{enumerate}
\item 범주 $\text{SR}(\mathcal{C})$는 쌍대곱을 가지며
$F$는 쌍대곱과 가환한다.
\item 함자 $F : \text{SR}(\mathcal{C}) \to \textit{PSh}(\mathcal{C})$는
극한과 가환한다.
\item $\mathcal{C}$가 섬유곱을 가지면 $\text{SR}(\mathcal{C})$도
섬유곱을 가진다.
\item $\mathcal{C}$가 이항곱을 가지면
$\text{SR}(\mathcal{C})$도 이항곱을 가진다.
\item $\mathcal{C}$가 등화자를 가지면 $\text{SR}(\mathcal{C})$도
등화자를 가진다.
\item $\mathcal{C}$가 끝 대상을 가지면 $\text{SR}(\mathcal{C})$도
끝 대상을 가진다.
\end{enumerate}
$X \in \Ob(\mathcal{C})$라 하자.
\begin{enumerate}
\item 범주 $\text{SR}(\mathcal{C}, X)$는 쌍대곱을 가지며
$F$는 쌍대곱과 가환한다.
\item $\mathcal{C}$가 섬유곱을 가지면 $\text{SR}(\mathcal{C}, X)$는
유한 극한을 가지며
$F : \text{SR}(\mathcal{C}, X) \to \textit{PSh}(\mathcal{C})/h_X$
는 이들과 가환한다.
\end{enumerate}
\end{lemma}

\begin{proof}
$\text{SR}(\mathcal{C})$에 관한 결과부터 증명한다.
(1)의 증명. $\{U_i\}_{i \in I}$와 $\{V_j\}_{j \in J}$의 쌍대곱은
$\{U_i\}_{i \in I} \amalg \{V_j\}_{j \in J}$이다. 다시 말해, 첨자 집합이
$I \amalg J$이고 원소 $k \in I \amalg J$에 대해 $k = i \in I$이면
$U_i$를, $k = j \in J$이면 $V_j$를 주는 대상들의 족이다.
대상들의 족으로 이루어진 족의 쌍대곱도 마찬가지이다.
$F$가 이들과 가환함은 명백하다.

\medskip\noindent
(2)의 증명. $U$가 $\Ob(\mathcal{C})$에 속할 때,
$\text{SR}(\mathcal{C})$의 대상 $\{U\}$를 생각하자. 분명히
$\Mor_{\text{SR}(\mathcal{C})}(\{U\}, K)) = F(K)(U)$
가 $K \in \Ob(\text{SR}(\mathcal{C}))$에 대해 성립한다.
준층의 극한은 단면 수준에서 계산되므로(Sites, 절
\ref{sites-section-limits-colimits-PSh}), $F$는 극한과 가환한다.

\medskip\noindent
(3)의 증명. 사상
$(\alpha, f_i) : \{U_i\}_{i \in I} \to \{V_j\}_{j \in J}$
와 사상
$(\beta, g_k) : \{W_k\}_{k \in K} \to \{V_j\}_{j \in J}$
가 주어졌다고 하자. 이 사상들의 섬유곱은
$$
\{ U_i \times_{f_i, V_j, g_k} W_k\}_{(i, j, k) \in I \times J \times K
\text{ such that } j = \alpha(i) = \beta(k)}
$$
로 주어진다. $\mathcal{C}$가 섬유곱을 가지면 이 섬유곱들이 존재한다.

\medskip\noindent
(4)의 증명. $\{U_i\}_{i \in I}$와 $\{V_j\}_{j \in J}$의 곱은
$\{U_i \times V_j\}_{i \in I, j \in J}$이다. $\mathcal{C}$가 곱을 가지면
이 곱들이 존재한다.

\medskip\noindent
(5)의 증명. 두 사상
$(\alpha, f_i), (\alpha', f'_i) : \{U_i\}_{i \in I} \to \{V_j\}_{j \in J}$
의 등화자는
$$
\{
\text{Eq}(f_i, f'_i : U_i \to V_{\alpha(i)})
\}_{i \in I,\ \alpha(i) = \alpha'(i)}
$$
$\mathcal{C}$가 등화자를 가지면 이 등화자들이 존재한다.

\medskip\noindent
(6)의 증명. $X$가 $\mathcal{C}$의 끝 대상이면 $\{X\}$는
$\text{SR}(\mathcal{C})$의 끝 대상이다.

\medskip\noindent
이제 $\text{SR}(\mathcal{C}, X)$에 관한 명제들을 증명한다.
위의 결과를 범주 $\mathcal{C}/X$에 적용하고
$\text{SR}(\mathcal{C}/X) = \text{SR}(\mathcal{C}, X)$ 및
$\textit{PSh}(\mathcal{C}/X) = \textit{PSh}(\mathcal{C})/h_X$
(혼돈 위상을 갖춘 $\mathcal{C}$에 Sites, 보조정리
\ref{sites-lemma-essential-image-j-shriek}를 적용한다)를 사용하면 따른다.
그러나 다음과 같이 직접 논증할 수도 있다. 분명히
$\{U_i \to X\}_{i \in I}$와 $\{V_j \to X\}_{j \in J}$
의 쌍대곱은
$\{U_i \to X\}_{i \in I} \amalg \{V_j \to X\}_{j \in J}$이고,
공역이 $X$인 사상들의 족으로 이루어진 족의 쌍대곱도 마찬가지이다.
대상 $\{X \to X\}$는 $\text{SR}(\mathcal{C}, X)$의 끝 대상이다.
사상
$(\alpha, f_i) : \{U_i \to X\}_{i \in I} \to \{V_j \to X\}_{j \in J}$
와 사상
$(\beta, g_k) : \{W_k \to X\}_{k \in K} \to \{V_j \to X\}_{j \in J}$
가 주어졌다고 하자. 이 사상들의 섬유곱은
$$
\{ U_i \times_{f_i, V_j, g_k} W_k \to X \}_{(i, j, k) \in I \times J \times K
\text{ such that } j = \alpha(i) = \beta(k)}
$$
로 주어진다. $\mathcal{C}$가 섬유곱을 가진다는 가정에 의해 이 섬유곱들이
존재한다. 따라서 $\text{SR}(\mathcal{C}, X)$는 유한 극한을 가진다.
Categories, 보조정리 \ref{categories-lemma-finite-limits-exist}를 보라.
이 경우 함자 $F$에 관한 명제의 확인은 생략한다.
\end{proof}





\section{초덮개}
\label{hypercovering-section-hypercoverings}

\noindent
우리 범주가 사이트라고 가정하면 다음 정의를 둘 수 있다.

\begin{definition}
\label{hypercovering-definition-covering-SR}
사이트 $\mathcal{C}$를 생각하자.
$f = (\alpha, f_i) : \{U_i\}_{i \in I} \to \{V_j\}_{j \in J}$
를 범주 $\text{SR}(\mathcal{C})$의 사상이라 하자.
모든 $j \in J$에 대해 사상의 족
$\{U_i \to V_j\}_{i \in I, \alpha(i) = j}$가 사이트
$\mathcal{C}$의 덮개이면 $f$를 {\it 덮개}라고 한다.
$X$를 $\mathcal{C}$의 대상이라 하자.
$\text{SR}(\mathcal{C}, X)$의 사상 $K \to L$의
$\text{SR}(\mathcal{C})$에서의 상이 덮개이면 이 사상을
{\it 덮개}라고 한다.
\end{definition}

\begin{lemma}
\label{hypercovering-lemma-covering-permanence}
사이트 $\mathcal{C}$를 생각하자.
\begin{enumerate}
\item $\text{SR}(\mathcal{C})$에서 덮개들의 합성은 덮개이다.
\item $K \to L$이 $\text{SR}(\mathcal{C})$의 덮개이고
$L' \to L$이 사상이면 $L' \times_L K$가 존재하며,
$L' \times_L K \to L'$은 덮개이다.
\item $\mathcal{C}$가 두 대상의 곱을 가지고
$A \to B$와 $K \to L$이 $\text{SR}(\mathcal{C})$의 덮개이면,
$A \times K \to B \times L$은 덮개이다.
\end{enumerate}
$X \in \Ob(\mathcal{C})$라 하자. 그러면 (1)과 (2)는
$\text{SR}(\mathcal{C}, X)$에 대해 성립하고, $\mathcal{C}$가
섬유곱을 가지면 (3)도 성립한다.
\end{lemma}

\begin{proof}
(1)은 사이트의 공리에서 곧바로 따른다.
(2)는 보조정리 \ref{hypercovering-lemma-coprod-prod-SR}의 증명에서 주어진
$\text{SR}(\mathcal{C})$의 섬유곱 구성과, $\mathcal{C}$의 덮개를
이루는 사상들이 표현 가능해야 한다는 조건에서 따른다.
$A \times K \to B \times L$을 합성
$A \times K \to B \times K \to B \times L$로 보면
덮개들의 밑변환을 합성한 것이므로 (3)이 따른다.
마지막 명제는 $\text{SR}(\mathcal{C}, X) =
\text{SR}(\mathcal{C}/X)$이기 때문에 성립한다.
\end{proof}

\noindent
보조정리 \ref{hypercovering-lemma-coprod-prod-SR}와
Simplicial, 보조정리 \ref{simplicial-lemma-existence-cosk}에 의해
$\mathcal{C}$가 섬유곱을 가지면 $\text{SR}(\mathcal{C}, X)$의
절단 단체적 대상의 여골격이 존재한다. 따라서 다음 정의는 의미를 가진다.

\begin{definition}
\label{hypercovering-definition-hypercovering}
$\mathcal{C}$를 사이트라 하자. $\mathcal{C}$가 섬유곱을 가진다고 가정하자.
$X \in \Ob(\mathcal{C})$를 $\mathcal{C}$의 대상이라 하자.
{\it $X$의 초덮개}란 다음을 만족하는
$\text{SR}(\mathcal{C}, X)$의 단체적 대상 $K$를 말한다.
\begin{enumerate}
\item 대상 $K_0$은 사이트 $\mathcal{C}$에서 $X$의 덮개이다.
\item 모든 $n \geq 0$에 대해 표준 사상
$$
K_{n + 1} \longrightarrow (\text{cosk}_n \text{sk}_n K)_{n + 1}
$$
은 위에서 정의한 의미의 덮개이다.
\end{enumerate}
\end{definition}

\noindent
$\text{SR}(\mathcal{C}, X)$의 각 대상은 결국 공역이 $X$인
사상들의 족이므로 조건 (1)은 의미를 가진다. 이 조건은 $K_0$에서
$\text{SR}(\mathcal{C}, X)$의 끝 대상으로 가는 사상이
덮개라고 표현할 수도 있다.

\begin{example}[체흐 초덮개]
\label{hypercovering-example-cech}
섬유곱을 갖는 사이트 $\mathcal{C}$를 생각하자.
$\{U_i \to X\}_{i \in I}$를 $\mathcal{C}$의 덮개라 하고
$K_0 = \{U_i \to X\}_{i \in I}$로 두자. 그러면 $K_0$은
$\text{SR}(\mathcal{C}, X)$의 $0$-절단 단체적 대상이므로
$$
K = \text{cosk}_0 K_0.
$$
를 만들 수 있다. $K$가 정의 \ref{hypercovering-definition-hypercovering}의 조건 (1)을
만족함은 분명하다. Simplicial, 보조정리
\ref{simplicial-lemma-cosk-up}에 의해 모든 사상
$K_{n + 1} \to (\text{cosk}_n \text{sk}_n K)_{n + 1}$이 동형이므로
조건 (2)도 만족한다. 각 항 $K_n$은 통상적인
$$
K_n = \{
U_{i_0} \times_X U_{i_1} \times_X \ldots \times_X U_{i_n} \to X
\}_{(i_0, i_1, \ldots, i_n) \in I^{n + 1}}
$$
임에 유의하자. 이 형태의 $X$의 초덮개를
$X$의 {\it 체흐 초덮개}라고 한다.
\end{example}

\begin{example}[사이트의 단체적 대상에 의한 초덮개]
\label{hypercovering-example-hypercovering-in-C}
섬유곱을 갖는 사이트 $\mathcal{C}$를 생각하자.
$X \in \Ob(\mathcal{C})$라 하자. $U$를 $\mathcal{C}$의 단체적 대상이라 하고,
통상적으로 $U_n = U([n])$라 쓰자. 마지막으로 증대
$$
a : U \to X
$$
가 주어졌다고 가정하자. 이 상황에서는 각 항이
$K_n = \{U_n \to X\}$인 $\text{SR}(\mathcal{C}, X)$의
단체적 대상 $K$를 생각할 수 있다. $K$가 정의
\ref{hypercovering-definition-hypercovering}의 의미에서 $X$의 초덮개일 필요충분조건은
다음 세 조건\footnote{$\mathcal{C}$가 섬유곱을 가지므로 범주
$\mathcal{C}/X$는 모든 유한 극한을 가진다. 따라서 Simplicial, 보조정리
\ref{simplicial-lemma-existence-cosk}에 의해 필요한 여골격들이 존재한다.}
이 성립하는 것이다.
\begin{enumerate}
\item $\{U_0 \to X\}$는 $\mathcal{C}$의 덮개이다.
\item $\{U_1 \to U_0 \times_X U_0\}$는 $\mathcal{C}$의 덮개이다.
\item $\{U_{n + 1} \to (\text{cosk}_n\text{sk}_n U)_{n + 1}\}$
는 $n \geq 1$에 대해 $\mathcal{C}$의 덮개이다.
\end{enumerate}
직접적인 확인은 생략한다.
\end{example}

\begin{example}[덮개 사상에 결부된 체흐 초덮개]
\label{hypercovering-example-cech-cover}
섬유곱을 갖는 사이트 $\mathcal{C}$를 생각하자. $U \to X$를
$\{U \to X\}$가 $\mathcal{C}$의 덮개가 되는 $\mathcal{C}$의 사상이라
하자\footnote{이 성질을 갖는 $\mathcal{C}$의 사상을 때로
``덮개 사상''이라고 한다.}. 각 항이
$$
K_n = \{U \times_X U \times_X \ldots \times_X U \to X\}
\quad (n + 1 \text{ factors})
$$
인 $\text{SR}(\mathcal{C}, X)$의 단체적 대상 $K$를 생각하자.
그러면 $K$는 $X$의 초덮개이다. 이 예는
예 \ref{hypercovering-example-cech}와 예 \ref{hypercovering-example-hypercovering-in-C} 모두의
특수한 경우이다.
\end{example}

\begin{lemma}
\label{hypercovering-lemma-hypercoverings-set}
섬유곱을 갖는 사이트 $\mathcal{C}$를 생각하자.
$X \in \Ob(\mathcal{C})$를 $\mathcal{C}$의 대상이라 하자.
$X$의 모든 초덮개로 이루어진 모임은 집합을 이룬다.
\end{lemma}

\begin{proof}
$\mathcal{C}$가 사이트이므로 $X$의 모든 덮개로 이루어진 모임은
집합을 이룬다. 따라서 가능한 $K_0$들의 모임은 집합이다.
가능한 모든 $K_0, \ldots, K_n$의 모임이 집합임을 보였다고 하자.
그러면 $K_0, \ldots, K_n$이 주어졌을 때 가능한 모든
$K_{n + 1}$의 모임이 집합임을 보이면 충분하다. 이는 분명하다.
실제로 $K_{n + 1}$은
$(\text{cosk}_n \text{sk}_n K)_{n + 1}$의 가능한 모든 덮개 중에서
택해야 한다.
\end{proof}

\begin{remark}
\label{hypercovering-remark-hypercoverings-really-set}
이 보조정리는 집합을 이루는 초덮개들의 선택으로 된 공종계가
존재한다고만 말하는 것이 아니라, 초덮개들 자체가 실제로
집합을 이룬다고 말한다.
\end{remark}

\noindent
$\mathcal{C}$ 위의 준층 범주는 유한 극한과 쌍대극한을 가진다.
따라서 집합 값 준층에 대해 함자 $\text{cosk}_n$이 존재한다.

\begin{lemma}
\label{hypercovering-lemma-hypercovering-F}
섬유곱을 갖는 사이트 $\mathcal{C}$를 생각하자.
$X \in \Ob(\mathcal{C})$를 $\mathcal{C}$의 대상이라 하고,
$K$를 $X$의 초덮개라 하자. 상수 단체적 준층 $h_X$로 가는
증대가 주어진 $\textit{PSh}(\mathcal{C})$의 단체적 대상
$F(K)$를 생각하자.
\begin{enumerate}
\item 준층의 사상 $F(K)_0 \to h_X$는 층화하면 전사가 된다.
\item 사상
$$
(d^1_0, d^1_1) :
F(K)_1
\longrightarrow
F(K)_0 \times_{h_X} F(K)_0
$$
은 층화하면 전사가 된다.
\item 모든 $n \geq 1$에 대해 사상
$$
F(K)_{n + 1} \longrightarrow (\text{cosk}_n \text{sk}_n F(K))_{n + 1}
$$
은 층화하면 전사가 된다.
\end{enumerate}
\end{lemma}

\begin{proof}
다음 사실을 사용한다. $\{U_i \to U\}_{i \in I}$가 사이트
$\mathcal{C}$의 덮개이면 사상
$$
\amalg_{i \in I} h_{U_i} \to h_U
$$
은 층화하면 전사가 된다. Sites, 보조정리
\ref{sites-lemma-covering-surjective-after-sheafification}를 보라.
따라서 첫 번째 명제는 곧바로 따른다.

\medskip\noindent
두 번째 명제에 대해서는 Simplicial, 예
\ref{simplicial-example-cosk0}에 따라 단체적 대상
$\text{cosk}_0 \text{sk}_0 K$의 각 항이
$K_0 \times \ldots \times K_0$임에 유의하자. 따라서 초덮개의 정의에 의해
$(d^1_0, d^1_1) : K_1 \to K_0 \times K_0$은 덮개이다.
위의 명제와 $F$가 곱을 $h_X$ 위의 섬유곱으로 보낸다는 사실에서
(2)가 따른다.

\medskip\noindent
세 번째 명제에 대해서는
$\text{cosk}_n \text{sk}_n F(K) =
F(\text{cosk}_n \text{sk}_n K)$가 $n \geq 1$에 대해 성립한다고
주장한다. 이를 증명하기 위해 함자
$\text{SR}(\mathcal{C}, X) \to \textit{PSh}(\mathcal{C})/h_X$를
잠시 $F'$로 나타내자. 보조정리 \ref{hypercovering-lemma-coprod-prod-SR}에 의해
함자 $F'$는 유한 극한과 가환한다. Simplicial, 절
\ref{simplicial-section-skeleton}에서 함자 $\text{cosk}_n$을
서술한 방식에 따라 $\text{cosk}_n \text{sk}_n F'(K) =
F'(\text{cosk}_n \text{sk}_n K)$임을 알 수 있다.
Simplicial, 보조정리 \ref{simplicial-lemma-existence-cosk}에서
$(\text{cosk}_n U)_m$을 서술할 때 쓰인 범주는
$(\Delta/[m])^{opp}_{\leq n}$임을 상기하자.
$n \geq 1$이면 $(\Delta/[m])_{\leq n}$이 연결 범주임을 보이는 것은
재미있는 연습이다(Categories, 정의
\ref{categories-definition-category-connected}를 보라).
따라서 Categories, 보조정리
\ref{categories-lemma-connected-limit-over-X}에 의해
$\text{cosk}_n \text{sk}_n F'(K) =
\text{cosk}_n \text{sk}_n F(K)$이고, 주장이 따른다.
이제 (2)의 사상은 정의 \ref{hypercovering-definition-covering-SR}에서 말하는
덮개에 함자 $F$를 적용하여 얻은 것이므로, 이로부터 성질 (2)가 따르고
결론은 이 증명의 첫머리에서 언급한 첫 번째 사실로부터 따른다.
\end{proof}



\section{비순환성}
\label{hypercovering-section-acyclicity}

\noindent
$\mathcal{C}$를 사이트라 하자. 집합의 준층 $\mathcal{F}$에 대해
다음 규칙으로 정의되는 아벨 군의 준층을
$\mathbf{Z}_\mathcal{F}$로 나타낸다.
$$
\mathbf{Z}_\mathcal{F}(U) = \text{free abelian group on }\mathcal{F}(U).
$$
이를 때로 {\it $\mathcal{F}$ 위의 자유 아벨 준층}이라고 한다.
물론 구성 $\mathcal{F} \mapsto \mathbf{Z}_\mathcal{F}$는 함자이며
망각 함자
$\textit{PAb}(\mathcal{C}) \to \textit{PSh}(\mathcal{C})$의
왼쪽 수반이다. 층화 $\mathbf{Z}_\mathcal{F}^\#$는 아벨 군의 층이고,
함자 $\mathcal{F} \mapsto \mathbf{Z}_\mathcal{F}^\#$ 역시 왼쪽 수반이다.
$\mathbf{Z}_\mathcal{F}^\#$를 때로
{\it $\mathcal{F}$ 위의 자유 아벨 층}이라고 한다.

\medskip\noindent
사이트 $\mathcal{C}$의 대상 $X$에 대해 $h_X$ 위의 자유 아벨 준층을
$\mathbf{Z}_X$로 나타내고, 그 층화를 $\mathbf{Z}_X^\#$로 나타낸다.

\begin{definition}
\label{hypercovering-definition-homology}
$\mathcal{C}$를 사이트라 하고, $K$를
$\textit{PSh}(\mathcal{C})$의 단체적 대상이라 하자.
위의 구성으로 $\textit{Ab}(\mathcal{C})$의 단체적 대상
$\mathbf{Z}_K^\#$를 얻는다. 이에 결부된 아벨 준층의 복합체
$s(\mathbf{Z}_K^\#)$를 취할 수 있다. Simplicial, 절
\ref{simplicial-section-complexes}를 보라. {\it $K$의 호몰로지}란
아벨 층의 복합체 $s(\mathbf{Z}_K^\#)$의 호몰로지이다.
\end{definition}

\noindent
다시 말해, {\it $K$의 $i$차 호몰로지 $H_i(K)$}란 아벨 군의 층
$H_i(K) = H_i(s(\mathbf{Z}_K^\#))$이다. 이 절에서는 $K$가
$\mathcal{C}$의 대상 $X$의 초덮개인 경우의 호몰로지를 다룬다.

\begin{lemma}
\label{hypercovering-lemma-compare-cosk0}
$\mathcal{C}$를 사이트라 하고, $\mathcal{F} \to \mathcal{G}$를
집합의 준층 사이의 사상이라 하자. $K$를
$\textit{PSh}(\mathcal{C})$의 단체적 대상 중에서 그 $n$차 항이
$\mathcal{F}$의 $\mathcal{G}$ 위 $(n + 1)$중 섬유곱인 것으로 나타내자.
Simplicial, 예 \ref{simplicial-example-fibre-products-simplicial-object}를
보라. $\mathcal{F} \to \mathcal{G}$가 층화하면 전사가 된다면
$$
H_i(K) =
\left\{
\begin{matrix}
0 & \text{if} & i > 0\\
\mathbf{Z}_\mathcal{G}^\# & \text{if} & i = 0
\end{matrix}
\right.
$$
이다. 차수 $0$에서의 동형은 사상
$H_0(K) \to \mathbf{Z}_\mathcal{G}^\#$로 주어지며, 이 사상은
사상 $(\mathbf{Z}_K^\#)_0 =
\mathbf{Z}_\mathcal{F}^\# \to \mathbf{Z}_\mathcal{G}^\#$에서 유도된다.
\end{lemma}

\begin{proof}
$\mathcal{G}' \subset \mathcal{G}$를 사상
$\mathcal{F} \to \mathcal{G}$의 상이라 하자.
$U \in \Ob(\mathcal{C})$를 택하고
$A = \mathcal{F}(U)$와 $B = \mathcal{G}'(U)$로 두자.
그러면 단체 집합 $K(U)$는 $n$-단체가
$$
A \times_B A \times_B \ldots \times_B A\ (n + 1 \text{ factors)}.
$$
로 주어지는 단체 집합과 같다. Simplicial, 보조정리
\ref{simplicial-lemma-cosk-minus-one-equivalence}에 의해
사상 $K(U) \to B$는 자명한 Kan 파이브레이션이다. 따라서 호모토피 동치이다
(Simplicial, 보조정리 \ref{simplicial-lemma-trivial-kan-homotopy}).
여기에 ``자유 아벨 군'' 함자를 적용하면
$$
\mathbf{Z}_K(U) \longrightarrow \mathbf{Z}_B
$$
도 호모토피 동치임을 얻는다. $s(\mathbf{Z}_B)$는 복합체
$$
\ldots \to
\bigoplus\nolimits_{b \in B}\mathbf{Z} \xrightarrow{0}
\bigoplus\nolimits_{b \in B}\mathbf{Z} \xrightarrow{1}
\bigoplus\nolimits_{b \in B}\mathbf{Z} \xrightarrow{0}
\bigoplus\nolimits_{b \in B}\mathbf{Z} \to 0
$$
임에 유의하자. Simplicial, 보조정리
\ref{simplicial-lemma-homology-eilenberg-maclane}를 보라.
따라서 $i > 0$에 대해 $H_i(s(\mathbf{Z}_K(U))) = 0$이고,
$H_0(s(\mathbf{Z}_K(U))) = \bigoplus_{b \in B}\mathbf{Z}
= \bigoplus_{s \in \mathcal{G}'(U)} \mathbf{Z}$.
이 식별들은 제한 사상과 호환된다.

\medskip\noindent
$i > 0$에 대해 $H_i(s(\mathbf{Z}_K)) = 0$이고
$H_0(s(\mathbf{Z}_K)) = \mathbf{Z}_{\mathcal{G}'}$라고 결론짓는다.
여기서는 $\textit{PAb}(\mathcal{C})$ 안에서 호몰로지군을 계산한다.
층화가 완전 함자이므로 보조정리의 결과가 따른다. 구체적으로 완전성은
$H_0(s(\mathbf{Z}_K))^\# = H_0(s(\mathbf{Z}_K^\#))$를 함의하며,
다른 지표도 마찬가지이다.
\end{proof}

\begin{lemma}
\label{hypercovering-lemma-acyclicity}
$\mathcal{C}$를 사이트라 하자. $f : L \to K$를
$\textit{PSh}(\mathcal{C})$의 단체적 대상 사이의 사상이라 하고,
$n \geq 0$을 정수라 하자. 다음을 가정하자.
\begin{enumerate}
\item $i < n$에 대해 사상 $L_i \to K_i$는 동형이다.
\item 사상 $L_n \to K_n$은 층화하면 전사가 된다.
\item 표준 사상 $L \to \text{cosk}_n \text{sk}_n L$은 동형이다.
\item 표준 사상 $K \to \text{cosk}_n \text{sk}_n K$은 동형이다.
\end{enumerate}
그러면 $H_i(f) : H_i(L) \to H_i(K)$는 동형이다.
\end{lemma}

\begin{proof}
이 증명은 위의 보조정리 \ref{hypercovering-lemma-compare-cosk0}의 증명과 완전히 같다.
먼저 $K_n' \subset K_n$을 사상 $L_n \to K_n$의 상인 부분 준층이라 하자.
가정 (2)는 $K_n'$의 층화가 $K_n$의 층화와 같다는 뜻이다.
또한 모든 $i < n$에 대해 $L_i = K_i$이므로
$U_0 = L_0 = K_0, \ldots, U_{n - 1} = L_{n - 1} = K_{n - 1}, U_n = K'_n$.
로 취하여 $n$-절단 단체적 준층 $U$를 얻는다.
단체적 준층 $K' = \text{cosk}_n U$로 두자. $K'_m$을 유한 극한으로
구성할 수 있고 층화는 완전하므로 $(K'_m)^\# = K_m$임을 알 수 있다.
다시 말해 $(K')^\# = K^\#$이다. 층화의 완전성으로 다시 한번
$H_i(K) = H_i(K')$라고 결론짓는다. 따라서 사상 $L \to K'$에 대해
보조정리를 증명하면 충분하다. 즉, $L_n \to K_n$이
{\it 준층}의 전사라고 가정해도 된다!

\medskip\noindent
이 경우 $\mathcal{C}$의 임의의 대상 $U$에 대해 단체 집합의 사상
$$
L(U) \longrightarrow K(U)
$$
은 Simplicial, 보조정리 \ref{simplicial-lemma-section}의
모든 가정을 만족한다. 따라서 자명한 Kan 파이브레이션이고, 특히
호모토피 동치이다(Simplicial, 보조정리
\ref{simplicial-lemma-trivial-kan-homotopy}). 그러므로
$$
\mathbf{Z}_L(U) \longrightarrow \mathbf{Z}_K(U)
$$
도 호모토피 동치이다. 이는 모든 $U$에 대해 성립하므로 결론이 따른다.
\end{proof}

\begin{lemma}
\label{hypercovering-lemma-acyclic-hypercover-sheaves}
$\mathcal{C}$를 사이트라 하고, $K$를 단체적 준층,
$\mathcal{G}$를 준층이라 하자. $K \to \mathcal{G}$를
$K$에서 $\mathcal{G}$로 가는 증대라 하자. 다음을 가정하자.
\begin{enumerate}
\item 준층의 사상 $K_0 \to \mathcal{G}$는 층화하면 전사가 된다.
\item 사상
$$
(d^1_0, d^1_1) :
K_1
\longrightarrow
K_0 \times_\mathcal{G} K_0
$$
은 층화하면 전사가 된다.
\item 모든 $n \geq 1$에 대해 사상
$$
K_{n + 1} \longrightarrow (\text{cosk}_n \text{sk}_n K)_{n + 1}
$$
은 층화하면 전사가 된다.
\end{enumerate}
그러면 $i > 0$에 대해 $H_i(K) = 0$이고
$H_0(K) = \mathbf{Z}_\mathcal{G}^\#$이다.
\end{lemma}

\begin{proof}
$n \geq 1$에 대해 $K^n = \text{cosk}_n \text{sk}_n K$로 나타내자.
$K^0$을 각 항 $(K^0)_n$이 $(n + 1)$중 섬유곱
$K_0 \times_\mathcal{G} \ldots \times_\mathcal{G} K_0$과 같은
단체적 대상으로 정의하자. Simplicial, 예
\ref{simplicial-example-fibre-products-simplicial-object}를 보라.
다음 사상들이 있다.
$$
K \longrightarrow \ldots \to K^n \to K^{n - 1} \to \ldots \to K^1 \to K^0.
$$
사상 $K \to K^i$와 $j \geq i \geq 1$일 때의 $K^j \to K^i$는
함자 $\text{cosk}_n$들의 보편 성질에서 유도된다.
사상 $K^1 \to K^0$은 Simplicial, 주석
\ref{simplicial-remark-augmentation}의 표준 사상이다.
또한 $K^0 \to \text{cosk}_1 \text{sk}_1 K^0$이 동형임을 상기하자.
Simplicial, 보조정리 \ref{simplicial-lemma-cosk-minus-one}를 보라.

\medskip\noindent
보조정리 \ref{hypercovering-lemma-compare-cosk0}에 의해 $i > 0$에 대해
$H_i(K^0) = 0$이고 $H_0(K^0) = \mathbf{Z}_\mathcal{G}^\#$임을 알 수 있다.

\medskip\noindent
$n \geq 1$을 택하고 사상 $K^n \to K^{n - 1}$을 생각하자.
이는 차수 $< n$인 항들에서 동형이다. 사상
$K^n \to \text{cosk}_n \text{sk}_n K^n$과
$K^{n - 1} \to \text{cosk}_n \text{sk}_n K^{n - 1}$
은 동형이다. 또한 $(K^n)_n = K_n$이고
$(K^{n - 1})_n = (\text{cosk}_{n - 1} \text{sk}_{n - 1} K)_n$이다.
따라서 가정에 의해 $(K^n)_n \to (K^{n - 1})_n$은
층화하면 전사가 되는 준층의 사상이다. 보조정리
\ref{hypercovering-lemma-acyclicity}에 의해
$H_i(K^n) = H_i(K^{n - 1})$라고 결론짓는다.
이를 위의 결과와 결합하면 보조정리가 증명된다.
\end{proof}

\begin{lemma}
\label{hypercovering-lemma-hypercovering-acyclic}
섬유곱을 갖는 사이트 $\mathcal{C}$를 생각하자.
$X$를 $\mathcal{C}$의 대상이라 하고, $K$를 $X$의 초덮개라 하자.
단체적 준층 $F(K)$의 호몰로지는 차수 $> 0$에서 $0$이고
차수 $0$에서는 $\mathbf{Z}_X^\#$과 같다.
\end{lemma}

\begin{proof}
보조정리 \ref{hypercovering-lemma-acyclic-hypercover-sheaves}와
\ref{hypercovering-lemma-hypercovering-F}를 결합하면 된다.
\end{proof}












\section{{\v C}ech 코호몰로지와 초덮개}
\label{hypercovering-section-hyper-cech}

\noindent
$\mathcal{C}$를 사이트라 하자. 사이트 $\mathcal{C}$ 위의 아벨 군 준층
$\mathcal{F}$를 생각하자. 이는 다음 함자를 정의한다.
\begin{eqnarray*}
\mathcal{F} : \text{SR}(\mathcal{C})^{opp}
& \longrightarrow &
\textit{Ab} \\
\{U_i\}_{i \in I} &
\longmapsto &
\prod\nolimits_{i \in I} \mathcal{F}(U_i)
\end{eqnarray*}
따라서 $\text{SR}(\mathcal{C})$의 단체적 대상 $K$는
$\textit{Ab}$의 코심플렉셜 대상 $\mathcal{F}(K)$로 옮겨진다.
$\mathcal{F}(K)$에 결부된 코체인 복합체 $s(\mathcal{F}(K))$를
(Simplicial, 절
\ref{simplicial-section-dold-kan-cosimplicial}) 단체적 대상 $K$에 관한
$\mathcal{F}$의 {\v C}ech 복합체라 한다. 다음과 같이 놓는다.
$$
\check{H}^i(K, \mathcal{F})
=
H^i(s(\mathcal{F}(K))).
$$
이를 $K$에 관한 $\mathcal{F}$의 $i$차 {\v C}ech 코호몰로지군이라 한다.
이 절에서는 Cohomology, 절 \ref{cohomology-section-cech},
\ref{cohomology-section-cech-functor},
\ref{cohomology-section-cech-cohomology-cohomology}에서 증명한 열린 덮개의
{\v C}ech 코호몰로지에 관한 몇몇 결과의 유사물을 증명한다.

\begin{lemma}
\label{hypercovering-lemma-h0-cech}
$\mathcal{C}$를 섬유곱을 갖는 사이트라 하자.
$X$를 $\mathcal{C}$의 대상이라 하자.
$K$를 $X$의 초덮개라 하자.
$\mathcal{F}$를 $\mathcal{C}$ 위의 아벨 군층이라 하자.
그러면 $\check{H}^0(K, \mathcal{F}) = \mathcal{F}(X)$이다.
\end{lemma}

\begin{proof}
다음이 성립한다.
$$
\check{H}^0(K, \mathcal{F})
=
\Ker(\mathcal{F}(K_0) \longrightarrow \mathcal{F}(K_1))
$$
$K_0 = \{U_i \to X\}$로 쓰자. 이는 사이트 $\mathcal{C}$의 덮개이다.
또한 $K_1 \to K_0 \times K_0$은
$\text{SR}(\mathcal{C}, X)$의 덮개이다. 따라서
$K_1 = \amalg_{i_0, i_1 \in I} \{V_{i_0i_1j} \to X\}$로 쓸 수 있고,
사상 $K_1 \to K_0 \times K_0$은 사이트 $\mathcal{C}$의 덮개들
$\{V_{i_0i_1j} \to U_{i_0} \times_X U_{i_1}\}$로 주어진다.
그러므로 다시 다음과 같이 식별할 수 있다.
$$
\check{H}^0(K, \mathcal{F})
=
\Ker(
\prod\nolimits_i \mathcal{F}(U_i)
\longrightarrow
\prod\nolimits_{i_0i_1 j} \mathcal{F}(V_{i_0i_1j})
)
$$
여기서 사상은 명백하다. $\mathcal{F}$의 층 조건은
$\check{H}^0(K, \mathcal{F}) = H^0(X, \mathcal{F})$를 함의한다.
\end{proof}

\noindent
실제로 이 성질은 물론 $\mathcal{C}$ 위의 모든 아벨 준층 가운데
아벨 층을 특징짓는다. 이 경우 Cohomology, 보조정리
\ref{hypercovering-lemma-injective-trivial-cech}의 유사물은 다음과 같다.

\begin{lemma}
\label{hypercovering-lemma-injective-trivial-cech}
$\mathcal{C}$를 섬유곱을 갖는 사이트라 하자.
$X$를 $\mathcal{C}$의 대상이라 하자.
$K$를 $X$의 초덮개라 하자.
$\mathcal{I}$를 $\mathcal{C}$ 위의 단사 아벨 군층이라 하자. 그러면
$$
\check{H}^p(K, \mathcal{I}) =
\left\{
\begin{matrix}
\mathcal{I}(X) & \text{if} & p = 0 \\
0 & \text{if} & p > 0
\end{matrix}
\right.
$$
\end{lemma}

\begin{proof}
$\text{SR}(\mathcal{C}, X)$의 임의의 대상 $Z = \{U_i \to X\}$와
$\mathcal{C}$ 위의 임의의 아벨 층 $\mathcal{F}$에 대해 다음이
성립함에 유의하자.
\begin{eqnarray*}
\mathcal{F}(Z)
& = &
\prod \mathcal{F}(U_i) \\
& = &
\prod \Mor_{\textit{PSh}(\mathcal{C})}(h_{U_i}, \mathcal{F})\\
& = &
\Mor_{\textit{PSh}(\mathcal{C})}(F(Z), \mathcal{F})\\
& = &
\Mor_{\textit{PAb}(\mathcal{C})}(\mathbf{Z}_{F(Z)}, \mathcal{F}) \\
& = &
\Mor_{\textit{Ab}(\mathcal{C})}(\mathbf{Z}_{F(Z)}^\#, \mathcal{F})
\end{eqnarray*}
따라서 $\text{SR}(\mathcal{C}, X)$의 임의의 단체적 대상 $K$에 대해
다음이 성립한다.
\begin{equation}
\label{hypercovering-equation-identify-cech}
s(\mathcal{F}(K))
=
\Hom_{\textit{Ab}(\mathcal{C})}(s(\mathbf{Z}_{F(K)}^\#), \mathcal{F})
\end{equation}
기호는 정의 \ref{hypercovering-definition-homology}을 보라.
$K$가 초덮개이면 층의 복합체 $s(\mathbf{Z}_{F(K)}^\#)$는
$\mathbf{Z}_X^\#$와 유사동형이다. 보조정리
\ref{hypercovering-lemma-hypercovering-acyclic}을 보라. 따라서 $\mathcal{I}$가 단사
아벨 층이고 $K$가 초덮개이면, 복합체 $s(\mathcal{I}(K))$는 차수 $0$을
제외하고 비순환이다. 다시 말해 다음이 성립한다.
$$
\check{H}^i(K, \mathcal{I}) = 0
$$
$i > 0$에 대해 그러하다. 이를 보조정리 \ref{hypercovering-lemma-h0-cech}와 결합하면
보조정리가 증명된다.
\end{proof}

\noindent
다음으로 Cohomology on Sites, 보조정리
\ref{sites-cohomology-lemma-cech-spectral-sequence}의 유사물을 다룬다.
$\mathcal{C}$를 사이트라 하자.
$\mathcal{F}$를 $\mathcal{C}$ 위의 아벨 군층이라 하자.
$\underline{H}^i(\mathcal{F})$는 규칙
$\underline{H}^i(\mathcal{F}) : U \longmapsto H^i(U, \mathcal{F})$로
정의되는 $\mathcal{C}$ 위의 아벨 군 준층을 나타냄을 상기하자.
이 절의 도입부에서와 같이 이를 $\text{SR}(\mathcal{C})$로 확장한다.

\begin{lemma}
\label{hypercovering-lemma-cech-spectral-sequence}
$\mathcal{C}$를 섬유곱을 갖는 사이트라 하자.
$X$를 $\mathcal{C}$의 대상이라 하자.
$K$를 $X$의 초덮개라 하자.
$\mathcal{F}$를 $\mathcal{C}$ 위의 아벨 군층이라 하자.
$D^{+}(\textit{Ab})$ 안에는 $\mathcal{F}$에 대해 함자적인 사상
$$
s(\mathcal{F}(K))
\longrightarrow
R\Gamma(X, \mathcal{F})
$$
이 있으며, 이는 함자 $\textit{Ab}(\mathcal{C}) \to \textit{Ab}$ 사이의
자연 변환
$$
\check{H}^i(K, -) \longrightarrow H^i(X, -)
$$
을 유도한다. 또한 다음을 만족하는 스펙트럼 열
$(E_r, d_r)_{r \geq 0}$이 존재한다.
$$
E_2^{p, q} = \check{H}^p(K, \underline{H}^q(\mathcal{F}))
$$
이는 $H^{p + q}(X, \mathcal{F})$로 수렴한다.
이 스펙트럼 열은 $\mathcal{F}$와 초덮개 $K$에 대해 함자적이다.
\end{lemma}

\begin{proof}
코호몰로지 장의 대응하는 보조정리에서 쓴 것과 같은 방법으로 증명할 수도
있다. 여기서는 대신 이중 복합체 논증으로 증명하자.

\medskip\noindent
$\mathcal{C}$ 위 아벨 층의 범주에서 단사 분해
$\mathcal{F} \to \mathcal{I}^\bullet$을 택하자. 항들이
$$
A^{p, q} = \mathcal{I}^q(K_p)
$$
인 이중 복합체 $A^{\bullet, \bullet}$를 생각하자. 여기서 미분
$d_1^{p, q} : A^{p, q} \to A^{p + 1, q}$는 코심플렉셜 아벨 군
$\mathcal{I}^p(K)$에 결부된 복합체 $s(\mathcal{I}^q(K))$의 미분에서 오고,
미분 $d_2^{p, q} : A^{p, q} \to A^{p, q + 1}$는 미분
$\mathcal{I}^q \to \mathcal{I}^{q + 1}$에서 온다.
$\text{Tot}(A^{\bullet, \bullet})$로 이중 복합체
$A^{\bullet, \bullet}$에 결부된 전체 복합체를 나타내자. Homology, 절
\ref{homology-section-double-complexes}을 보라. 이 이중 복합체에 결부된
두 스펙트럼 열 $({}'E_r, {}'d_r)$과 $({}''E_r, {}''d_r)$을 사용할 것이다.
Homology, 절 \ref{homology-section-double-complex}을 보라.

\medskip\noindent
보조정리 \ref{hypercovering-lemma-injective-trivial-cech}에 의해 복합체
$s(\mathcal{I}^q(K))$는 양의 차수에서 비순환이고, 그 $H^0$는
$\mathcal{I}^q(X)$와 같다. 따라서 Homology, 보조정리
\ref{homology-lemma-double-complex-gives-resolution}에 의해 자연 사상
$$
\mathcal{I}^\bullet(X) \longrightarrow \text{Tot}(A^{\bullet, \bullet})
$$
은 아벨 군 복합체의 유사동형이다. 특히
$H^n(\text{Tot}(A^{\bullet, \bullet})) = H^n(X, \mathcal{F})$임을 얻는다.

\medskip\noindent
보조정리의 사상 $s(\mathcal{F}(K)) \longrightarrow R\Gamma(X, \mathcal{F})$는
사상 $s(\mathcal{F}(K)) \to \text{Tot}(A^{\bullet, \bullet})$에 이어서
위에 표시한 유사동형의 역을 합성한 것이다. 이는
$\mathcal{I}^\bullet(X)$가 $R\Gamma(X, \mathcal{F})$의 대표이기 때문이다.

\medskip\noindent
스펙트럼 열 $({}'E_r, {}'d_r)_{r \geq 0}$을 생각하자. Homology,
보조정리 \ref{homology-lemma-ss-double-complex}에 의해 다음이 성립한다.
$$
{}'E_2^{p, q} = H^p_I(H^q_{II}(A^{\bullet, \bullet}))
$$
다시 말해 먼저 $d_2$에 관해 코호몰로지를 취하면 군
${}'E_1^{p, q} = \underline{H}^q(\mathcal{F})(K_p)$를 얻는다.
따라서 미분 ${}'d_1$의 서술에 의해 실제로
${}'E_2^{p, q} = \check{H}^p(K, \underline{H}^q(\mathcal{F}))$이다.
위의 결과와 Homology, 보조정리
\ref{homology-lemma-first-quadrant-ss}에 의해 이는 원하는 대로
$H^n(X, \mathcal{F})$로 수렴한다.

\medskip\noindent
위 구성들이 아벨 층 $\mathcal{F}$와 초덮개 $K$에 대해 함자적이라는
명제의 증명은 생략한다.
\end{proof}









\section{베르디에식 초덮개}
\label{hypercovering-section-hypercoverings-verdier}

\noindent
통찰력 있는 독자라면 대상 $X$의 초덮개에 대한 체흐-코호몰로지에서
코호몰로지로 가는 스펙트럼 열을 얻는 데 필요한 것은 보조정리
\ref{hypercovering-lemma-hypercovering-F}의 결론뿐임을 알아차렸을 것이다.
따라서 다음 정의는 의미가 있다.

\begin{definition}
\label{hypercovering-definition-hypercovering-variant}
$\mathcal{C}$를 사이트라 하자. $\mathcal{C}$가 등화자와 섬유곱을
갖는다고 가정하자. $\mathcal{G}$를 집합의 준층이라 하자.
{\it $\mathcal{G}$의 초덮개}란 다음 조건을 만족하는 증대
$F(K) \to \mathcal{G}$가 주어진 $\text{SR}(\mathcal{C})$의 단체적 대상
$K$를 말한다.
\begin{enumerate}
\item $F(K_0) \to \mathcal{G}$는 층화한 뒤 전사가 된다.
\item $F(K_1) \to F(K_0) \times_\mathcal{G} F(K_0)$
는 층화한 뒤 전사가 된다.
\item $F(K_{n + 1}) \longrightarrow F((\text{cosk}_n \text{sk}_n K)_{n + 1})$
는 $n \geq 1$에 대해 층화한 뒤 전사가 된다.
\end{enumerate}
$\text{SR}(\mathcal{C})$의 단체적 대상 $K$가
$\textit{PSh}(\mathcal{C})$의 끝대상 $*$의 초덮개이면 $K$를
{\it 초덮개}라 한다.
\end{definition}

\noindent
$\mathcal{C}$가 섬유곱과 등화자를 갖는다는 가정은
$\text{SR}(\mathcal{C})$도 섬유곱과 등화자를 가지며 $F$가 이들과
가환함을 보장한다(보조정리 \ref{hypercovering-lemma-coprod-prod-SR}). 이는 사용하는
여골격 함자들을 정의하기에 충분하다(Simplicial, 주석
\ref{simplicial-remark-existence-cosk}과 Categories, 보조정리
\ref{categories-lemma-fibre-products-equalizers-exist}을 보라).
$\mathcal{C}$가 일반적인 경우에는 조건 (3)을 다음 조건으로 바꾸어도
이 절의 결과들이 그대로 성립한다.
$F(K_{n + 1}) \longrightarrow ((\text{cosk}_n \text{sk}_n F(K))_{n + 1})$
는 $n \geq 1$에 대해 층화한 뒤 전사가 된다.

\medskip\noindent
$\mathcal{F}$를 $\mathcal{C}$ 위의 아벨 층이라 하자.
앞 절에서 $\text{SR}(\mathcal{C})$의 단체적 대상 $K$에 관한
$\mathcal{F}$의 {\v C}ech 복합체를 정의했다. 이제 준층 $\mathcal{G}$가
주어졌을 때 다음과 같이 놓는다.
$$
H^0(\mathcal{G}, \mathcal{F}) =
\Mor_{\textit{PSh}(\mathcal{C})}(\mathcal{G}, \mathcal{F}) =
\Mor_{\Sh(\mathcal{C})}(\mathcal{G}^\#, \mathcal{F}) =
H^0(\mathcal{G}^\#, \mathcal{F})
$$
기호는 Cohomology on Sites, 절
\ref{sites-cohomology-section-limp}에서와 같다. 이는 왼쪽 완전 함자이고,
그 고차 유도 함자들(Cohomology on Sites, 절
\ref{sites-cohomology-section-limp}에서 간단히 연구했다)은
$H^i(\mathcal{G}, \mathcal{F})$로 나타낸다. $\mathcal{G}$의 초덮개 $K$가
주어지면 코호몰로지 $H^i(\mathcal{G}, \mathcal{F})$로 수렴하는
체흐-코호몰로지에서 코호몰로지로 가는 스펙트럼 열이 있음을 보일 것이다.
$\mathcal{G} = *$이면
$H^i(*, \mathcal{F}) = H^i(\mathcal{C}, \mathcal{F})$가 사이트
$\mathcal{C}$ 위에서 $\mathcal{F}$의 코호몰로지를 되찾음에 유의하자.

\begin{lemma}
\label{hypercovering-lemma-h0-cech-variant}
$\mathcal{C}$를 등화자와 섬유곱을 갖는 사이트라 하자.
$\mathcal{G}$를 $\mathcal{C}$ 위의 준층이라 하자.
$K$를 $\mathcal{G}$의 초덮개라 하자.
$\mathcal{F}$를 $\mathcal{C}$ 위의 아벨 군층이라 하자.
그러면 $\check{H}^0(K, \mathcal{F}) = H^0(\mathcal{G}, \mathcal{F})$이다.
\end{lemma}

\begin{proof}
$H^0(\mathcal{G}, \mathcal{F})$의 정의와 다음 도식이 층화한 뒤
쌍대등화자 도식이 된다는 사실에서 따른다.
$$
\xymatrix{
F(K_1) \ar@<1ex>[r] \ar@<-1ex>[r] &
F(K_0) \ar[r] & \mathcal{G}
}
$$
\end{proof}

\begin{lemma}
\label{hypercovering-lemma-injective-trivial-cech-variant}
$\mathcal{C}$를 등화자와 섬유곱을 갖는 사이트라 하자.
$\mathcal{G}$를 $\mathcal{C}$ 위의 준층이라 하자.
$K$를 $\mathcal{G}$의 초덮개라 하자.
$\mathcal{I}$를 $\mathcal{C}$ 위의 단사 아벨 군층이라 하자. 그러면
$$
\check{H}^p(K, \mathcal{I}) =
\left\{
\begin{matrix}
H^0(\mathcal{G}, \mathcal{I}) & \text{if} & p = 0 \\
0 & \text{if} & p > 0
\end{matrix}
\right.
$$
\end{lemma}

\begin{proof}
(\ref{hypercovering-equation-identify-cech})에 의해 다음이 성립한다.
$$
s(\mathcal{F}(K))
=
\Hom_{\textit{Ab}(\mathcal{C})}(s(\mathbf{Z}_{F(K)}^\#), \mathcal{F})
$$
복합체 $s(\mathbf{Z}_{F(K)}^\#)$는 $\mathbf{Z}_\mathcal{G}^\#$와
유사동형이다. 보조정리 \ref{hypercovering-lemma-acyclic-hypercover-sheaves}를 보라.
따라서 $\mathcal{I}$가 단사 아벨 층이면 복합체 $s(\mathcal{I}(K))$는
차수 $0$을 제외하고 비순환이다. 다시 말해 $i > 0$에 대해
$\check{H}^i(K, \mathcal{I}) = 0$이다. 이를 보조정리
\ref{hypercovering-lemma-h0-cech-variant}와 결합하면 보조정리가 증명된다.
\end{proof}

\begin{lemma}
\label{hypercovering-lemma-cech-spectral-sequence-variant}
$\mathcal{C}$를 등화자와 섬유곱을 갖는 사이트라 하자.
$\mathcal{G}$를 $\mathcal{C}$ 위의 준층이라 하자.
$K$를 $\mathcal{G}$의 초덮개라 하자.
$\mathcal{F}$를 $\mathcal{C}$ 위의 아벨 군층이라 하자.
$D^{+}(\textit{Ab})$ 안에는 $\mathcal{F}$에 대해 함자적인 사상
$$
s(\mathcal{F}(K)) \longrightarrow R\Gamma(\mathcal{G}, \mathcal{F})
$$
이 있으며, 이는 함자 $\textit{Ab}(\mathcal{C}) \to \textit{Ab}$ 사이의
자연 변환
$$
\check{H}^i(K, -) \longrightarrow H^i(\mathcal{G}, -)
$$
을 유도한다. 또한 다음을 만족하는 스펙트럼 열
$(E_r, d_r)_{r \geq 0}$이 존재한다.
$$
E_2^{p, q} = \check{H}^p(K, \underline{H}^q(\mathcal{F}))
$$
이는 $H^{p + q}(\mathcal{G}, \mathcal{F})$로 수렴한다.
이 스펙트럼 열은 $\mathcal{F}$와 초덮개 $K$에 대해 함자적이다.
\end{lemma}

\begin{proof}
$\mathcal{C}$ 위 아벨 층의 범주에서 단사 분해
$\mathcal{F} \to \mathcal{I}^\bullet$을 택하자. 항들이
$$
A^{p, q} = \mathcal{I}^q(K_p)
$$
인 이중 복합체 $A^{\bullet, \bullet}$를 생각하자. 여기서 미분
$d_1^{p, q} : A^{p, q} \to A^{p + 1, q}$는 미분
$\mathcal{I}^p \to \mathcal{I}^{p + 1}$에서 오고, 미분
$d_2^{p, q} : A^{p, q} \to A^{p, q + 1}$는 위에서 설명했듯이
코심플렉셜 아벨 군 $\mathcal{I}^p(K)$에 결부된 복합체
$s(\mathcal{I}^p(K))$의 미분에서 온다. 이 이중 복합체에 결부된
두 스펙트럼 열 $({}'E_r, {}'d_r)$과 $({}''E_r, {}''d_r)$을 사용할 것이다.
Homology, 절 \ref{homology-section-double-complex}을 보라.

\medskip\noindent
보조정리 \ref{hypercovering-lemma-injective-trivial-cech-variant}에 의해 복합체
$s(\mathcal{I}^p(K))$는 양의 차수에서 비순환이고, 그 $H^0$는
$H^0(\mathcal{G}, \mathcal{I}^p)$와 같다. 따라서 Homology, 보조정리
\ref{homology-lemma-double-complex-gives-resolution}와 그 증명에 의해
스펙트럼 열 $({}'E_r, {}'d_r)$은 퇴화하며, 자연 사상
$$
H^0(\mathcal{G}, \mathcal{I}^\bullet) \longrightarrow
\text{Tot}(A^{\bullet, \bullet})
$$
은 아벨 군 복합체의 유사동형이다. 보조정리의 사상
$s(\mathcal{F}(K)) \longrightarrow R\Gamma(\mathcal{G}, \mathcal{F})$는
자연 사상 $s(\mathcal{F}(K)) \to \text{Tot}(A^{\bullet, \bullet})$에
이어서 위에 표시한 유사동형의 역을 합성한 것이다. 이는
$H^0(\mathcal{G}, \mathcal{I}^\bullet)$이
$R\Gamma(\mathcal{G}, \mathcal{F})$의 대표이기 때문이다.

\medskip\noindent
스펙트럼 열 $({}''E_r, {}''d_r)_{r \geq 0}$을 생각하자. Homology,
보조정리 \ref{homology-lemma-ss-double-complex}에 의해 다음이 성립한다.
$$
{}''E_2^{p, q} = H^p_{II}(H^q_I(A^{\bullet, \bullet}))
$$
다시 말해 먼저 $d_1$에 관해 코호몰로지를 취하면 군
${}''E_1^{p, q} = \underline{H}^p(\mathcal{F})(K_q)$를 얻는다.
따라서 미분 ${}''d_1$의 서술에 의해 실제로
${}''E_2^{p, q} = \check{H}^p(K, \underline{H}^q(\mathcal{F}))$이다.
이 스펙트럼 열은 $\text{Tot}(A^{\bullet, \bullet})$의 코호몰로지로
수렴하므로 증명이 끝난다.
\end{proof}

\begin{lemma}
\label{hypercovering-lemma-cech-spectral-sequence-verdier}
$\mathcal{C}$를 등화자와 섬유곱을 갖는 사이트라 하자.
$K$를 초덮개라 하자.
$\mathcal{F}$를 아벨 층이라 하자. 다음을 만족하는 스펙트럼 열
$(E_r, d_r)_{r \geq 0}$이 존재한다.
$$
E_2^{p, q} = \check{H}^p(K, \underline{H}^q(\mathcal{F}))
$$
이는 대역 코호몰로지군 $H^{p + q}(\mathcal{F})$로 수렴한다.
\end{lemma}

\begin{proof}
이는 보조정리 \ref{hypercovering-lemma-cech-spectral-sequence-variant}의 특수한 경우이다.
\end{proof}







\section{덮개에 의한 초덮개의 세분}
\label{hypercovering-section-covering}

\noindent
초덮개를 구성하는 몇 가지 방법을 제시한다. 범주
$\text{SR}(\mathcal{C}, X)$가 섬유곱을 가지므로
$\text{SR}(\mathcal{C}, X)$의 단체적 대상들의 범주도
섬유곱을 가진다는 점에 유의하자. Simplicial, 보조정리
\ref{simplicial-lemma-fibre-product}을 보라.

\begin{lemma}
\label{hypercovering-lemma-funny-gamma}
$\mathcal{C}$를 섬유곱을 가지는 사이트라 하고,
$X$를 $\mathcal{C}$의 대상이라 하자.
$K, L, M$을 $\text{SR}(\mathcal{C}, X)$의 단체적 대상들이라 하고,
$a : K \to L$, $b : M \to L$을 사상들이라 하자. 다음을 가정하자.
\begin{enumerate}
\item $K$는 $X$의 초덮개이다.
\item 사상 $M_0 \to L_0$은 덮개이다.
\item 모든 $n \geq 0$에 대해 다음 도식에서
$$
\xymatrix{
M_{n + 1} \ar[dd] \ar[rr] \ar[rd]^\gamma &
&
(\text{cosk}_n \text{sk}_n M)_{n + 1} \ar[dd] \\
&
L_{n + 1}
\times_{(\text{cosk}_n \text{sk}_n L)_{n + 1}}
(\text{cosk}_n \text{sk}_n M)_{n + 1}
\ar[ld] \ar[ru]
& \\
L_{n + 1} \ar[rr] & & (\text{cosk}_n \text{sk}_n L)_{n + 1}
}
$$
화살표 $\gamma$는 덮개이다.
\end{enumerate}
그러면 섬유곱 $K \times_L M$은 $X$의 초덮개이다.
\end{lemma}

\begin{proof}
(2)에 의해 사상 $(K \times_L M)_0 = K_0 \times_{L_0} M_0 \to K_0$은
덮개의 밑변환이므로 덮개이다. 보조정리
\ref{hypercovering-lemma-covering-permanence}을 보라. 또한 (1)에 의해
$K_0 \to \{X \to X\}$는 덮개이다. 따라서 보조정리
\ref{hypercovering-lemma-covering-permanence}에 의해
$(K \times_L M)_0 \to \{X \to X\}$는 덮개이다. 그러므로
$K \times_L M$은 정의 \ref{hypercovering-definition-hypercovering}의 첫 번째 조건을
만족한다.

\medskip\noindent
이제 모든 $n \geq 0$에 대해
$$
K_{n + 1} \times_{L_{n + 1}} M_{n + 1} = (K \times_L M)_{n + 1}
\longrightarrow
(\text{cosk}_n \text{sk}_n (K \times_L M))_{n + 1}
$$
가 덮개임을 확인해야 한다. 다음과 같이 줄여 쓰자.
$A = (\text{cosk}_n \text{sk}_n K)_{n + 1}$,
$B = (\text{cosk}_n \text{sk}_n L)_{n + 1}$, 그리고
$C = (\text{cosk}_n \text{sk}_n M)_{n + 1}$.
함자 $\text{cosk}_n \text{sk}_n$은 섬유곱과 가환한다.
Simplicial, 보조정리 \ref{simplicial-lemma-cosk-fibre-product}을 보라.
따라서 위의 오른쪽 항은 $A \times_B C$와 같다. 다음 가환 도식을
생각하자.
$$
\xymatrix{
K_{n + 1} \times_{L_{n + 1}} M_{n + 1} \ar[r] \ar[d] &
M_{n + 1} \ar[d] \ar[rd]_\gamma \ar[rrd] &
& \\
K_{n + 1} \ar[r] \ar[rd] &
L_{n + 1} \ar[rrd] &
L_{n + 1} \times_B C \ar[l] \ar[r] &
C \ar[d] \\
&
A \ar[rr] &
&
B
}
$$
이 도식은
$$
K_{n + 1} \times_{L_{n + 1}} M_{n + 1}
=
(K_{n + 1} \times_B C)
\times_{(L_{n + 1} \times_B C), \gamma}
M_{n + 1}
$$
임을 보여 준다. 이제 $K_{n + 1} \times_B C \to A \times_B C$는
덮개 $K_{n + 1} \to A$의 사상 $A \times_B C \to A$에 의한
밑변환이므로 덮개이다. 가정 (3)에 의해 사상 $\gamma$는 덮개이다.
따라서 사상
$$
(K_{n + 1} \times_B C)
\times_{(L_{n + 1} \times_B C), \gamma}
M_{n + 1}
\longrightarrow
K_{n + 1} \times_B C
$$
은 덮개의 밑변환이므로 덮개이다. 덮개들의 합성은 덮개이므로
보조정리가 따른다.
\end{proof}

\begin{lemma}
\label{hypercovering-lemma-product-hypercoverings}
$\mathcal{C}$를 섬유곱을 가지는 사이트라 하고,
$X$를 $\mathcal{C}$의 대상이라 하자. $K, L$이 $X$의 초덮개들이면
$K \times L$은 $X$의 초덮개이다.
\end{lemma}

\begin{proof}
직접 확인해도 되고, 위의 보조정리 \ref{hypercovering-lemma-funny-gamma}를 사용하여
$L \to \{X \to X\}$가 성질 (3)을 가짐을 확인해도 된다.
\end{proof}


\noindent
$\mathcal{C}$를 섬유곱을 가지는 사이트라 하고,
$X$를 $\mathcal{C}$의 대상이라 하자. 범주
$\text{SR}(\mathcal{C}, X)$가 쌍대곱과 유한 극한을 가지므로,
어떤 단체 집합 $U$(예를 들어 비퇴화 단체를 유한 개만 가지는 것)와
$\text{SR}(\mathcal{C}, X)$의 임의의 단체적 대상 $K$에 대해 대상
$U \times K$와 $\Hom(U, K)$을 말할 수 있다. Simplicial, 절
\ref{simplicial-section-product-with-simplicial-sets}와
\ref{simplicial-section-hom-from-simplicial-sets}을 보라.

\begin{lemma}
\label{hypercovering-lemma-covering}
$\mathcal{C}$를 섬유곱을 가지는 사이트라 하고,
$X$를 $\mathcal{C}$의 대상이라 하자. $K$를 $X$의 초덮개라 하고,
$k \geq 0$을 정수라 하자. $u : Z \to K_k$를
$\text{SR}(\mathcal{C}, X)$의 덮개라 하자. 그러면
$L_k \to K_k$가 $u$를 거쳐 인수분해되는 초덮개들의 사상
$f: L \to K$가 존재한다.
\end{lemma}

\begin{proof}
$Y = K_k$로 나타내자. $C[k]$를 Simplicial, 예
\ref{simplicial-example-simplex-cosimplicial-set}에서 정의한
코심플렉셜 집합이라 하자. Simplicial, 보조정리
\ref{simplicial-lemma-morphism-into-product}에 주어진
$\Hom(C[k], Y)$와 $\Hom(C[k], Z)$의 기술을 사용하겠다.
$\text{id} : K_k = Y \to Y$에 대응하는 정준 사상
$K \to \Hom(C[k], Y)$가 있다. 사상
$\Hom(C[k], Z) \to \Hom(C[k], Y)$을 생각하자. 그 $n$차 항에서의 사상은
$$
\prod\nolimits_{\alpha : [k] \to [n]} Z
\longrightarrow
\prod\nolimits_{\alpha : [k] \to [n]} Y
$$
이며, 각 인자에서 주어진 사상 $Z \to Y$를 사용한다. 다음과 같이 두자.
$$
L = K \times_{\Hom(C[k], Y)} \Hom(C[k], Z).
$$
사상 $L_k \to K_k$는 다음 가환 도식에 놓인다.
$$
\xymatrix{
L_k \ar[r] \ar[d] &
\prod_{\alpha : [k] \to [k]} Z \ar[r]^-{\text{pr}_{\text{id}_{[k]}}} \ar[d] &
Z \ar[d] \\
K_k \ar[r] &
\prod_{\alpha : [k] \to [k]} Y \ar[r]^-{\text{pr}_{\text{id}_{[k]}}} &
Y
}
$$
아래쪽 두 화살표의 합성이 항등사상이므로, 원하는 인수분해를 얻는다.

\medskip\noindent
이제 $L$이 $X$의 초덮개임을 보여야 한다. 이를 위해 보조정리
\ref{hypercovering-lemma-funny-gamma}를 사용하겠다. 조건 (1)은 가정에 의해 만족된다.
(2)에 대해서는 사상
$$
\Hom(C[k], Z)_0 \to \Hom(C[k], Y)_0
$$
이 덮개이다. 실제로 사상 $[k] \to [0]$은 하나뿐이므로 이는
$Z \to Y$와 동형이다.

\medskip\noindent
이제 $n = 0$에 대한 조건 (3)을 생각하자. 이때
$(\text{cosk}_0 T)_1 = T \times T$
(Simplicial, 예 \ref{simplicial-example-cosk0})이고
$\Hom(C[k], Z)_1 = \prod_{\alpha : [k] \to [1]} Z$이므로 다음 도식을 얻는다.
$$
\xymatrix{
\prod\nolimits_{\alpha : [k] \to [1]} Z \ar[r] \ar[d] &
Z \times Z \ar[d] \\
\prod\nolimits_{\alpha : [k] \to [1]} Y \ar[r] &
Y \times Y
}
$$
여기서 가로 화살표들은 두 비전사 $\alpha$에 대응하는 인자들로의
사영에 대응한다. 따라서 화살표 $\gamma$는 사상
$$
\prod\nolimits_{\alpha : [k] \to [1]} Z
\longrightarrow
\prod\nolimits_{\alpha : [k] \to [1]\text{ not onto}} Z
\times
\prod\nolimits_{\alpha : [k] \to [1]\text{ onto}} Y
$$
이며, 이는 덮개들의 곱이므로 보조정리
\ref{hypercovering-lemma-covering-permanence}에 의해 덮개이다.

\medskip\noindent
이제 $n > 0$에 대한 조건 (3)을 생각하자. 유한 집합들의 단사 사상
$\tau : S' \to S$로서, $\text{SR}(\mathcal{C}, X)$의 임의의 대상
$T$에 대해 사상
\begin{equation}
\label{hypercovering-equation-map}
\Hom(C[k], T)_{n + 1} \to
(\text{cosk}_n \text{sk}_n \Hom(C[k], T))_{n + 1}
\end{equation}
이 $T$에 대해 함자적으로 사영
$\prod_{s \in S} T \to \prod_{s' \in S'} T$와 동형이 되게 하는 것이
존재한다고 주장한다. 이것이 참이면 앞 문단과 같은 논증으로 화살표
$\gamma$가 사상
$$
\prod\nolimits_{s \in S} Z
\longrightarrow
\prod\nolimits_{s \in S'} Z
\times
\prod\nolimits_{s \not\in \tau(S')} Y
$$
임을 알 수 있다. 이는 덮개들의 곱이므로 보조정리
\ref{hypercovering-lemma-covering-permanence}에 의해 덮개이다. 구성에 의해
$\Hom(C[k], T)_{n + 1} = \prod_{\alpha : [k] \to [n + 1]} T$
(Simplicial, 보조정리 \ref{simplicial-lemma-morphism-into-product}을 보라).
따라서 $S = \text{Map}([k], [n + 1])$로 취한다. 한편 Simplicial,
보조정리 \ref{simplicial-lemma-formula-limit}은
$(\text{cosk}_n \text{sk}_n \Hom(C[k], T))_{n + 1}$
의 점들을 $\Hom(C[k], T)_n$의 점들의 열
$(f_0, \ldots, f_{n + 1})$로 기술한다. 여기서
$0 \leq i < j \leq n + 1$에 대해
$d^n_{j - 1} f_i = d^n_i f_j$이다. $f_i = (f_{i, \alpha})$로 쓸 수 있는데,
여기서 $f_{i, \alpha}$는 $T$의 점이고
$\alpha \in \text{Map}([k], [n])$이다. 이 조건들은
$$
f_{i, \delta^n_{j - 1} \circ \beta} = f_{j, \delta_i^n \circ \beta}
$$
가 임의의 $0 \leq i < j \leq n + 1$과
$\beta : [k] \to [n - 1]$에 대해 성립한다는 조건으로 바뀐다. 따라서
$$
S' = \{0, \ldots, n + 1\} \times \text{Map}([k], [n]) / \sim
$$
이며, 여기서 동치 관계는 다음 동치들로 생성된다.
$$
(i, \delta^n_{j - 1} \circ \beta) \sim (j, \delta_i^n \circ \beta)
$$
여기서 $0 \leq i < j \leq n + 1$이고 $\beta : [k] \to [n - 1]$이다.
계산(생략됨)에 의하면 사상 (\ref{hypercovering-equation-map})은 $(i, \alpha)$를
$\delta^{n + 1}_i \circ \alpha \in S$로 보내는 사상 $S' \to S$에
대응한다. (Simplicial, 보조정리
\ref{simplicial-lemma-relations-face-degeneracy}의 (1)에 의해 이 사상이
잘 정의됨을 확인할 수 있으므로 독자는 안심할 수 있을 것이다.)
증명을 끝내기 위해서는
$\alpha, \alpha' : [k] \to [n]$와 $0 \leq i < j \leq n + 1$
가
$$
\delta^{n + 1}_i \circ \alpha = \delta^{n + 1}_j \circ \alpha'
$$
을 만족하면 어떤 $\beta : [k] \to [n - 1]$에 대해
$\alpha = \delta^n_{j - 1} \circ \beta$이고
$\alpha' = \delta_i^n \circ \beta$임을 보이면 충분하다.
이는 쉽게 확인할 수 있으므로 생략한다.
\end{proof}

\begin{lemma}
\label{hypercovering-lemma-covering-sheaf}
$\mathcal{C}$를 섬유곱을 가지는 사이트라 하고,
$X$를 $\mathcal{C}$의 대상이라 하자. $K$를 $X$의 초덮개라 하고,
$n \geq 0$을 정수라 하자. $u : \mathcal{F} \to F(K_n)$를
층화하면 전사가 되는 준층들의 사상이라 하자. 그러면
$F(f_n) : F(L_n) \to F(K_n)$이 $u$를 거쳐 인수분해되는
초덮개들의 사상 $f: L \to K$가 존재한다.
\end{lemma}

\begin{proof}
$K_n = \{U_i \to X\}_{i \in I}$로 쓰자. 그러면 사상 $u$는
집합의 준층들의 사상 $u : \mathcal{F} \to \amalg h_{u_i}$이다.
$u$에 관한 가정은 모든 $i \in I$에 대해 사이트 $\mathcal{C}$의 덮개
$\{U_{ij} \to U_i\}_{j \in I_i}$와 준층들의 사상
$t_{ij} : h_{U_{ij}} \to \mathcal{F}$가 존재하여,
$u \circ t_{ij}$가 사상 $U_{ij} \to U_i$에서 오는 사상
$h_{U_{ij}} \to h_{U_i}$가 된다는 뜻이다.
$J = \amalg_{i \in I} I_i$로 두고, $\alpha : J \to I$를 자명한
사상이라 하자. $j \in J$에 대해 $V_j = U_{\alpha(j)j}$로 나타내고,
$Z = \{V_j \to X\}_{j \in J}$로 두자. 마지막으로
$\alpha : J \to I$와 위의 사상들
$V_j = U_{\alpha(j)j} \to U_{\alpha(j)}$로 주어지는 사상
$u' : Z \to K_n$을 생각하자. 이것은 분명 범주
$\text{SR}(\mathcal{C}, X)$의 덮개이고, 구성에 의해
$F(u') : F(Z) \to F(K_n)$은 $u$를 거쳐 인수분해된다.
따라서 결과는 위의 보조정리 \ref{hypercovering-lemma-covering}에서 따른다.
\end{proof}


\section{단체의 첨가}
\label{hypercovering-section-adding-simplices}

\noindent
이 절에서는 뒤에 필요할 몇 가지 기술적인 보조정리를 증명한다.
$\mathcal{C}$를 섬유곱을 가지는 사이트라 하고,
$X$를 $\mathcal{C}$의 대상이라 하자. 위의 절
\ref{hypercovering-section-covering}에서 지적했듯이, 어떤 단체 집합 $U$와
$\text{SR}(\mathcal{C}, X)$의 임의의 단체적 대상 $K$에 대해 대상
$U \times K$와 $\Hom(U, K)$가 정의된다. Simplicial, 절
\ref{simplicial-section-product-with-simplicial-sets}와
\ref{simplicial-section-hom-from-simplicial-sets}을 보라.

\begin{lemma}
\label{hypercovering-lemma-one-more-simplex}
$\mathcal{C}$를 섬유곱을 가지는 사이트라 하고,
$X$를 $\mathcal{C}$의 대상이라 하자. $K$를 $X$의 초덮개라 하자.
$U \subset V$를 단체 집합들이라 하고, 모든 $n$에 대해
$U_n, V_n$이 유한하고 공집합이 아니라고 하자.
$U$가 비퇴화 단체를 유한 개만 가진다고 가정하자.
$n \geq 0$, $x \in V_n$, $x \not \in U_n$이 다음을 만족한다고 하자.
\begin{enumerate}
\item $i < n$에 대해 $V_i = U_i$이다.
\item $V_n = U_n \cup \{x\}$이다.
\item $j > n$이고 $z \in V_j$, $z \not \in U_j$이면 그 원소는 퇴화되어 있다.
\end{enumerate}
그러면 $\text{SR}(\mathcal{C}, X)$의 사상
$$
\Hom(V, K)_0
\longrightarrow
\Hom(U, K)_0
$$
은 덮개이다.
\end{lemma}

\begin{proof}
$n = 0$이면 $V = U \amalg \Delta[0]$임이 쉽게 따른다(아래를 보라).
이 경우 $\Hom(V, K)_0 =
\Hom(U, K)_0 \times K_0$이다. 따라서 이 경우의 결과는 보조정리
\ref{hypercovering-lemma-covering-permanence}에서 따른다.

\medskip\noindent
$a : \Delta[n] \to V$를 Simplicial, 보조정리
\ref{simplicial-lemma-simplex-map}에서와 같이 $x$에 결부된 사상이라 하자.
$\partial \Delta[n] = i_{(n-1)!} \text{sk}_{n - 1} \Delta[n]$로 쓰고,
이를 $\Delta[n]$의 $(n - 1)$-골격이라 하자.
$b : \partial \Delta[n] \to U$를 $a$의 $\Delta[n]$의
$(n - 1)$-골격으로의 제한이라 하자. Simplicial, 보조정리
\ref{simplicial-lemma-glue-simplex}에 의해
$V = U \amalg_{\partial \Delta[n]} \Delta[n]$이다. Simplicial,
보조정리 \ref{simplicial-lemma-hom-from-coprod}에 의해
$$
\xymatrix{
\Hom(V, K)_0 \ar[r] \ar[d] &
\Hom(U, K)_0 \ar[d] \\
\Hom(\Delta[n], K)_0 \ar[r] &
\Hom(\partial \Delta[n], K)_0
}
$$
은 섬유곱 정사각형이다. 따라서 아래쪽 가로 화살표가 덮개임을
보이면 충분하다. Simplicial, 보조정리
\ref{simplicial-lemma-cosk-shriek}에 의해 이 화살표는
$$
K_n \to (\text{cosk}_{n - 1} \text{sk}_{n - 1} K)_n
$$
와 동일시되므로, 초덮개의 정의에 의해 덮개이다.
\end{proof}

\begin{lemma}
\label{hypercovering-lemma-add-simplices}
$\mathcal{C}$를 섬유곱을 가지는 사이트라 하고,
$X$를 $\mathcal{C}$의 대상이라 하자. $K$를 $X$의 초덮개라 하자.
$U \subset V$를 단체 집합들이라 하고, 모든 $n$에 대해
$U_n, V_n$이 유한하고 공집합이 아니라고 하자.
$U$와 $V$가 비퇴화 단체를 유한 개만 가진다고 가정하자.
그러면 $\text{SR}(\mathcal{C}, X)$의 사상
$$
\Hom(V, K)_0
\longrightarrow
\Hom(U, K)_0
$$
은 덮개이다.
\end{lemma}

\begin{proof}
위의 보조정리 \ref{hypercovering-lemma-one-more-simplex}에 의해, 이 보조정리에서와
같은 단체 집합들의 포함 $U \subset V$에 관한 간단한 보조정리를
증명하면 충분하다. 이것은 바로 Simplicial, 보조정리
\ref{simplicial-lemma-add-simplices}의 결과이다.
\end{proof}

\begin{lemma}
\label{hypercovering-lemma-degeneracy-maps-coverings}
$\mathcal{C}$를 섬유곱을 가지는 사이트라 하고, $X$를
$\mathcal{C}$의 대상이라 하자. $K$를 $X$의 초덮개라 하자. 그러면
\begin{enumerate}
\item 각 $n \geq 0$에 대해 $K_n$은 $X$의 덮개이다.
\item 모든 $n \geq 1$과 $0 \leq i \leq n$에 대해
$d^n_i : K_n \to K_{n - 1}$은 덮개이다.
\end{enumerate}
\end{lemma}

\begin{proof}
$K_0$은 정의 \ref{hypercovering-definition-hypercovering}에 의해 $X$의 덮개임을
상기하자. 이는 $K_0 \to \{X \to X\}$가 정의
\ref{hypercovering-definition-covering-SR}의 의미에서 덮개라는 것과 동치이다.
따라서 (1)은 (2)에서 따른다. 실제로 (2)는 합성
$K_n \to K_{n - 1} \to \ldots \to K_0 \to \{X \to X\}$
이 보조정리 \ref{hypercovering-lemma-covering-permanence}에 의해 덮개임을 보인다.

\medskip\noindent
(2)의 증명. Simplicial, 보조정리
\ref{simplicial-lemma-exists-hom-from-simplicial-set-finite}에 의해
$\Mor(\Delta[n], K)_0 = K_n$임에 유의하자. 따라서 보조정리
\ref{hypercovering-lemma-add-simplices}를 $n + 1$개의 서로 다른 포함
$\Delta[n - 1] \to \Delta[n]$에 적용하면 (2)가 따른다.
\end{proof}

\begin{remark}
\label{hypercovering-remark-P-covering}
보조정리 \ref{hypercovering-lemma-add-simplices}와
\ref{hypercovering-lemma-degeneracy-maps-coverings}의 유용한 특수한 경우는 다음과 같다.
$\mathcal{C}$를 섬유곱을 가지는 범주라 하자.
$P \subset \text{Arrows}(\mathcal{C})$를 밑변환과 합성에 대해
안정적이고 모든 동형사상을 포함하는 부분집합이라 하자. 이때
{\it $P$-초덮개}란 $\mathcal{C}$의 단체적 대상에서 오는 증대
$a : U \to X$로서 다음을 만족하는 것을 말한다.
\begin{enumerate}
\item $U_0 \to X$는 $P$에 속한다.
\item $U_1 \to U_0 \times_X U_0$은 $P$에 속한다.
\item $U_{n + 1} \to (\text{cosk}_n\text{sk}_n U)_{n + 1}$
은 $n \geq 1$에 대해 $P$에 속한다.
\end{enumerate}
범주 $\mathcal{C}/X$는 모든 유한 극한을 가지므로 위의 정식화에서
사용한 여골격들이 존재한다(Categories, 보조정리
\ref{categories-lemma-finite-limits-exist}을 보라).
이제 사상들 $U_n \to X$와 $d^n_i : U_n \to U_{n - 1}$이 $P$에
속한다고 주장한다. $\mathcal{C}$를 $f \in P$인
$\{f : V \to U\}$들을 덮개로 가지는 사이트로 만들고,
$K_n = \{U_n \to X\}$로 주어지는 $K$를 취하면, 이 주장은 앞서 말한
보조정리들에서 따른다.
\end{remark}


\section{호모토피}
\label{hypercovering-section-homotopies}

\noindent
$\mathcal{C}$를 섬유곱을 가지는 사이트라 하고,
$X$를 $\mathcal{C}$의 대상이라 하자. $L$을
$\text{SR}(\mathcal{C}, X)$의 단체적 대상이라 하자.
Simplicial, 보조정리
\ref{simplicial-lemma-exists-hom-from-simplicial-set-finite}에 의해 함자
$$
T
\longmapsto
\Mor_{\text{Simp}(\text{SR}(\mathcal{C}, X))}(\Delta[1] \times T, L)
$$
를 표현하는 대상 $\Hom(\Delta[1], L)$이 범주
$\text{Simp}(\text{SR}(\mathcal{C}, X))$에 존재한다.
$e_i : \Delta[0] \to \Delta[1]$와 동일시
$\Hom(\Delta[0], L) = L$에서 오는 정준 사상
$$
\Hom(\Delta[1], L) \to L \times L
$$
이 있다.

\begin{lemma}
\label{hypercovering-lemma-hom-hypercovering}
$\mathcal{C}$를 섬유곱을 가지는 사이트라 하고,
$X$를 $\mathcal{C}$의 대상이라 하자. $L$을
$\text{SR}(\mathcal{C}, X)$의 단체적 대상이라 하고,
$n \geq 0$이라 하자. 위에서 정의한 사상으로부터 오는 가환 도식
\begin{equation}
\label{hypercovering-equation-diagram}
\xymatrix{
\Hom(\Delta[1], L)_{n + 1} \ar[r] \ar[d] &
(\text{cosk}_n \text{sk}_n \Hom(\Delta[1], L))_{n + 1} \ar[d] \\
(L \times L)_{n + 1} \ar[r] &
(\text{cosk}_n \text{sk}_n (L \times L))_{n + 1}
}
\end{equation}
을 생각하자. 이 도식의 항들은 다음과 같이 동일시할 수 있다. 여기서
$\partial \Delta[n + 1] = i_{n!}\text{sk}_n \Delta[n + 1]$
은 $(n + 1)$-단체의 $n$-골격이다.
\begin{eqnarray*}
\Hom(\Delta[1], L)_{n + 1}
& = &
\Hom(\Delta[1] \times \Delta[n + 1], L)_0 \\
(\text{cosk}_n \text{sk}_n \Hom(\Delta[1], L))_{n + 1}
& = &
\Hom(\Delta[1] \times \partial \Delta[n + 1], L)_0 \\
(L \times L)_{n + 1}
& = &
\Hom(
(\Delta[n + 1] \amalg \Delta[n + 1], L)_0 \\
(\text{cosk}_n \text{sk}_n (L \times L))_{n + 1}
& = &
\Hom(
\partial \Delta[n + 1]
\amalg
\partial \Delta[n + 1], L)_0
\end{eqnarray*}
또한 $\text{SR}(\mathcal{C}, X)$의 이 대상들 사이의 사상들은
단체 집합들의 가환 도식
\begin{equation}
\label{hypercovering-equation-dual-diagram}
\xymatrix{
\Delta[1] \times \Delta[n + 1] &
\Delta[1] \times \partial\Delta[n + 1] \ar[l] \\
\Delta[n + 1] \amalg \Delta[n + 1] \ar[u] &
\partial\Delta[n + 1] \amalg \partial\Delta[n + 1]
\ar[l] \ar[u]
}
\end{equation}
에서 온다. 더욱이 도식 (\ref{hypercovering-equation-diagram})의 아래쪽 화살표와
오른쪽 화살표의 섬유곱은
$$
\Hom(U, L)_0
$$
와 같다. 여기서 $U \subset \Delta[1] \times \Delta[n + 1]$는
$\Delta[n + 1] \amalg \Delta[n + 1]$와
$\Delta[1] \times \partial\Delta[n + 1]$가 모두 그 안으로 사상되는
가장 작은 단체적 부분집합이다.
\end{lemma}

\begin{proof}
첫 번째와 세 번째 등식은 Simplicial, 보조정리
\ref{simplicial-lemma-exists-hom-from-simplicial-set-finite}이다.
두 번째와 네 번째 등식은 인용한 보조정리와 Simplicial, 보조정리
\ref{simplicial-lemma-cosk-shriek}를 함께 적용하면 따른다.
마지막 주장은 $U$가 도식 (\ref{hypercovering-equation-dual-diagram})의 아래쪽
화살표와 오른쪽 화살표의 밀어내기라는 사실에서 Simplicial, 보조정리
\ref{simplicial-lemma-hom-from-coprod}를 통해 따른다. $U$가 이
밀어내기와 같음을 보이려면,
$\Delta[n + 1] \amalg \Delta[n + 1]$와
$\Delta[1] \times \partial\Delta[n + 1]$
의 $\Delta[1] \times \Delta[n + 1]$ 안에서의 교집합이
$\partial\Delta[n + 1] \amalg \partial\Delta[n + 1]$.
와 같음을 보면 충분하다. 이는 독자에게 맡긴다.
\end{proof}

\begin{lemma}
\label{hypercovering-lemma-homotopy}
$\mathcal{C}$를 섬유곱을 가지는 사이트라 하고,
$X$를 $\mathcal{C}$의 대상이라 하자. $K, L$을 $X$의 초덮개들이라 하고,
$a, b : K \to L$을 초덮개들의 사상들이라 하자. 그러면
$a \circ c$가 $b \circ c$와 호모토픽이 되게 하는 초덮개들의 사상
$c : K' \to K$가 존재한다.
\end{lemma}

\begin{proof}
다음 가환 도식을 생각하자.
$$
\xymatrix{
K' \ar@{=}[r]^-{def} \ar[rd]_c &
K \times_{(L \times L)} \Hom(\Delta[1], L)
\ar[r] \ar[d] & \Hom(\Delta[1], L) \ar[d] \\
& K \ar[r]^{(a, b)} & L \times L
}
$$
$\Hom(\Delta[1], L)$의 함자성에 의해 가로 사상들의 합성은 사상
$K' \times \Delta[1] \to L$에 대응하며, 이 사상은
$c \circ a$와 $c \circ b$ 사이의 호모토피를 정의한다. 따라서
$K'$이 $X$의 초덮개임을 보이면 보조정리를 얻는다.
이를 위해 보조정리 \ref{hypercovering-lemma-funny-gamma}를 사상들의 쌍
$K \to L \times L$와 $\Hom(\Delta[1], L) \to L \times L$에 적용하겠다.
보조정리 \ref{hypercovering-lemma-funny-gamma}의 조건 (1)은 만족된다.
보조정리 \ref{hypercovering-lemma-funny-gamma}의 조건 (2)도 참이다. 실제로
$\Hom(\Delta[1], L)_0 = L_1$이고, $L$이
초덮개라는 가정에 의해 사상
$(d^1_0, d^1_1) : L_1 \to L_0 \times L_0$
은 $\text{SR}(\mathcal{C}, X)$의 덮개이기 때문이다.
보조정리 \ref{hypercovering-lemma-funny-gamma}의 조건 (3)을 증명하기 위해 위의
보조정리 \ref{hypercovering-lemma-hom-hypercovering}을 사용한다. 이 보조정리에 따르면,
보조정리 \ref{hypercovering-lemma-funny-gamma}의 조건 (3)에 나오는 사상 $\gamma$는
$$
\Hom(\Delta[1] \times \Delta[n + 1], L)_0
\longrightarrow
\Hom(U, L)_0
$$
이다. 여기서 $U \subset \Delta[1] \times \Delta[n + 1]$이다.
보조정리 \ref{hypercovering-lemma-add-simplices}에 의해 이는 덮개이므로 주장이
증명되었다.
\end{proof}

\begin{remark}
\label{hypercovering-remark-contractible-category}
증명의 핵심은 보조정리 \ref{hypercovering-lemma-add-simplices}를 사용하는 데
있음에 유의하자. 이 보조정리는 완전히 일반적이며, 단체 집합들이
비퇴화 단체를 유한 개만 가지는 한 그 정확한 모양과 무관하다.
따라서 다음과 같은 결과를 기대하는 것은 매우 합리적이다.
$K$와 $L$이 초덮개일 때 임의의 사상
$a : K \times \partial \Delta[k] \to L$이 주어졌다고 하자.
그러면 초덮개들의 사상 $c : K' \to K$와 사상
$g : K' \times \Delta[k] \to L$가 존재하고, 다음 등식이 성립한다.
$g|_{K' \times \partial \Delta[k]} =
a \circ (c \times \text{id}_{\partial \Delta[k]})$.
다시 말해 초덮개들의 범주는 적절한 의미에서 축약 가능하다.
\end{remark}

















\section{코호몰로지와 초덮개}
\label{hypercovering-section-cohomology}

\noindent
$\mathcal{C}$를 섬유곱을 갖는 사이트라 하자.
$X$를 $\mathcal{C}$의 대상이라 하고,
$\mathcal{F}$를 $\mathcal{C}$ 위의 아벨 군의 층이라 하자.
$K, L$을 $X$의 초덮개라 하자.
$a, b : K \to L$이 호모토픽 사상이면,
$\mathcal{F}(a), \mathcal{F}(b) : \mathcal{F}(K) \to \mathcal{F}(L)$도
호모토픽 사상이다. Simplicial, Lemma
\ref{simplicial-lemma-functorial-homotopy}를 보라.
따라서 결부된 코체인 복합체들의 코호몰로지 군에 같은 작용을 한다.
Simplicial, Lemma \ref{simplicial-lemma-homotopy-s-Q}를 보라.
이를 사용하여 모든 초덮개에 걸친 여극한을 정의할 것이다.

\medskip\noindent
$X$의 초덮개들을 대상으로 하고, $X$의 초덮개 사이의 사상들을
호모토피로 나눈 것을 사상으로 하는 범주를 잠시
$\text{HC}(\mathcal{C}, X)$라 쓰자. 이것이 범주이며 ``큰'' 범주가
아님을 이미 보았다. Lemma \ref{hypercovering-lemma-hypercoverings-set}를 보라.
$\text{HC}(\mathcal{C}, X)$의 반대 범주를 다음 도식의 첨자 범주로
사용한다. 용어에 대해서는 Categories, Section
\ref{categories-section-limits}를 보라. 도식
$$
\check{H}^i(-, \mathcal{F}) :
\text{HC}(\mathcal{C}, X)^{opp}
\longrightarrow
\textit{Ab}.
$$
을 생각하자. Lemmas \ref{hypercovering-lemma-product-hypercoverings}와
\ref{hypercovering-lemma-homotopy}, 그리고 위의 호모토피에 관한 설명에 의해
이 도식은 유향이다. Categories, Definition
\ref{categories-definition-directed}를 보라. 따라서 여극한
$$
\check{H}^i_{\text{HC}}(X, \mathcal{F}) =
\colim_{K \in \text{HC}(\mathcal{C}, X)} \check{H}^i(K, \mathcal{F})
$$
은 특히 간단히 기술된다(인용된 곳을 보라).

\begin{theorem}
\label{hypercovering-theorem-cohomology-hypercoverings}
$\mathcal{C}$를 섬유곱을 갖는 사이트라 하자.
$X$를 $\mathcal{C}$의 대상이라 하고 $i \geq 0$이라 하자. 함자들
\begin{eqnarray*}
\textit{Ab}(\mathcal{C}) & \longrightarrow & \textit{Ab} \\
\mathcal{F} & \longmapsto & H^i(X, \mathcal{F}) \\
\mathcal{F} & \longmapsto & \check{H}^i_{\text{HC}}(X, \mathcal{F})
\end{eqnarray*}
은 표준 동형이다.
\end{theorem}

\begin{proof}[스펙트럼 열을 사용하는 증명]
어떤 $p \geq 0$에 대해 $\xi \in H^p(X, \mathcal{F})$라 하자.
$X$의 어떤 초덮개 $K$에 대해 $\xi$가 Lemma
\ref{hypercovering-lemma-cech-spectral-sequence}의 사상
$\check{H}^p(X, \mathcal{F}) \to H^p(X, \mathcal{F})$의 상에
속함을 보이자.

\medskip\noindent
$p = 0$이면 Lemma \ref{hypercovering-lemma-h0-cech}에 의해 성립한다.
$p = 1$이면 $\xi|_{U_i} = 0$을 만족하는 사이트 $\mathcal{C}$의 덮개
$K_0 = \{U_i \to X\}$에서 시작하여 Example \ref{hypercovering-example-cech}에서와
같은 $X$의 {\v C}ech 초덮개 $K$를 택한다. Cohomology on Sites,
Lemma \ref{sites-cohomology-lemma-kill-cohomology-class-on-covering}를 보라.
이 경우 Lemma \ref{hypercovering-lemma-cech-spectral-sequence}의 스펙트럼 열에서
$\xi$가 $\check{H}^1(K, \mathcal{F})$의 한 원소에서 온다는 것이 곧바로
따른다. 일반적으로는 $\xi$가
$\underline{H}^p(\mathcal{F})(K_0)$에서 0으로 가도록 하는 $X$의
임의의 초덮개 $K$를 택한다(여기서도 Example \ref{hypercovering-example-cech}와
Cohomology on Sites, Lemma
\ref{sites-cohomology-lemma-kill-cohomology-class-on-covering}를 사용한다).
Lemma \ref{hypercovering-lemma-cech-spectral-sequence}의 스펙트럼 열에 따르면,
$\xi$가 $\check{H}^p(K, \mathcal{F})$의 한 원소에서 오는 데 대한
장애물은 다음을 만족하는 원소들의 열
$\xi_1, \ldots, \xi_{p - 1}$이다.
$\xi_q \in \check{H}^{p - q}(K, \underline{H}^q(\mathcal{F}))$
(더 정확히는 이 군들의 어떤 부분몫에서 $\xi_q$의 상들이다.)

\medskip\noindent
초덮개 $K$를 세분들로 귀납적으로 바꾸어 장애물
$\xi_1, \ldots, \xi_{p - 1}$의 제한이 0이 되게 할 수 있다
(부분몫에서의 상들만 0이 되는 것이 아니므로 여기에는 미묘한 점이 없다).
실제로 이미 $\xi_{q + 1}, \ldots, \xi_{p - 1}$이 0인 상황에
도달했다고 하자. $\xi_q \in
\check{H}^{p - q}(K, \underline{H}^q(\mathcal{F}))$는 어떤 원소
$$
\tilde \xi_q \in
\underline{H}^q(\mathcal{F})(K_{p - q}) =
\prod H^q(U_i, \mathcal{F})
$$
의 류임에 유의하자. 여기서
$K_{p - q} = \{U_i \to X\}_{i \in I}$이다. $\xi_{q, i}$를
$\tilde \xi_q$의 $H^q(U_i, \mathcal{F})$에서의 성분이라 하자.
$q \geq 1$이므로 다시 Cohomology on Sites, Lemma
\ref{sites-cohomology-lemma-kill-cohomology-class-on-covering}를 사용하여,
각 제한 $\xi_{q, i}|_{U_{i, j}} = 0$이 되도록 하는 사이트의 덮개들
$\{U_{i, j} \to U_i\}$를 택할 수 있다. 범주
$\text{SR}(\mathcal{C}, X)$의 대상 $Z = \{U_{i, j} \to X\}$와 그 명백한 사상
$u : Z \to K_{p - q}$를 생각하자. Definition
\ref{hypercovering-definition-covering-SR}에 의해 $u$는 분명히 덮개이다. Lemma
\ref{hypercovering-lemma-covering}에 의해 $L_{p - q} \to K_{p - q}$가 $u$를 거쳐
인수분해되는 $X$의 초덮개 사상 $L \to K$가 존재한다. 그러면
$\xi_q$의 $\underline{H}^q(\mathcal{F})(L_{p - q})$에서의 상은
분명히 0이다. Lemma \ref{hypercovering-lemma-cech-spectral-sequence}의 스펙트럼 열은
함자적이므로, $K$를 $L$로 바꾼 뒤에는
$\xi_q, \ldots, \xi_{p - 1}$이 모두 0인 상황에 도달한다.
이 과정을 계속하면 마침내 이들이 모두 0인 초덮개를 얻고, 따라서
$\xi$는 사상 $\check{H}^p(X, \mathcal{F}) \to H^p(X, \mathcal{F})$의
상에 속한다.

\medskip\noindent
$K$가 $X$의 초덮개이고 $\xi \in \check{H}^p(K, \mathcal{F})$이며,
Lemma \ref{hypercovering-lemma-cech-spectral-sequence}의 사상
$\check{H}^p(X, \mathcal{F}) \to H^p(X, \mathcal{F})$ 아래에서
$\xi$의 상이 0이라고 하자. 정리의 증명을 마치려면
$\xi$의 $\check{H}^p(L, \mathcal{F})$로의 제한이 0이 되게 하는
초덮개의 사상 $L \to K$가 존재함을 보여야 한다. Lemma
\ref{hypercovering-lemma-cech-spectral-sequence}의 스펙트럼 열에 따르면,
$\xi$의 $H^p(X, \mathcal{F})$에서의 상이 소멸한다는 것은 다음을
만족하는 원소들 $\xi_1, \ldots, \xi_{p - 2}$가 존재한다는 뜻이다.
$\xi_q \in \check{H}^{p - 1 - q}(K, \underline{H}^q(\mathcal{F}))$
(더 정확히는 이들의 어떤 부분몫에서의 상들)이고, 스펙트럼 열에서
상들 $d_{q + 1}^{p - 1 - q, q}\xi_q$의 합이 $\xi$이다. 따라서 위와
완전히 같은 방법으로, 원소들 $\xi_q$, $q = 1, \ldots, p - 2$의
$\check{H}^{p - 1 - q}(L, \underline{H}^q(\mathcal{F}))$에서의 제한이
0이 되게 하는 초덮개의 사상 $L \to K$를 찾을 수 있다. Lemma
\ref{hypercovering-lemma-cech-spectral-sequence}에 따라 사상 $L \to K$가
스펙트럼 열의 사상을 유도하므로, 이제 $\xi$는 0이다.
\end{proof}

\begin{proof}[스펙트럼 열을 사용하지 않는 증명]
$i = 0$인 경우의 결과는 이미 보았다. Lemma \ref{hypercovering-lemma-h0-cech}를 보라.
함자들 $H^i(X, -)$가 보편 $\delta$-함자를 이룬다는 것을 알고 있다.
Derived Categories, Lemma \ref{derived-lemma-higher-derived-functors}를 보라.
정리를 증명하려면 함자들의 열 $\check{H}^i_{HC}(X, -)$가
$\delta$-함자를 이룬다는 것을 보이면 충분하다. 실제로 {\v C}ech
코호몰로지가 단사층에서 0임을 알고 있으므로(Lemma
\ref{hypercovering-lemma-injective-trivial-cech}), Homology, Lemma
\ref{homology-lemma-efface-implies-universal}를 적용할 수 있다.

\medskip\noindent
다음을 $\mathcal{C}$ 위의 아벨 층들의 짧은 완전열이라 하자.
$$
0 \to \mathcal{F} \to \mathcal{G} \to \mathcal{H} \to 0
$$
$\xi \in \check{H}^p_{HC}(X, \mathcal{H})$라 하자. $X$의 초덮개
$K$와 코호몰로지에서 $\xi$를 나타내는 원소
$\sigma \in \mathcal{H}(K_p)$를 택하자. 이에 대응하는 복합체의 완전열
$$
0 \to s(\mathcal{F}(K)) \to s(\mathcal{G}(K)) \to s(\mathcal{H}(K))
$$
이 있지만, 오른쪽에도 0이 있다고 보장되지는 않는다. 이것만이 뱀
보조정리를 단순히 적용하여 $\delta(\xi)$를 정의하는 일을 막는다.
$K_p = \{U_i \to X\}$이면
$$
\mathcal{H}(K_p) = \prod \mathcal{H}(U_i)
$$
임을 상기하자. $\sigma =\prod \sigma_i$라 쓰자. 여기서
$\sigma_i \in \mathcal{H}(U_i)$이다. $\mathcal{G} \to \mathcal{H}$가
층의 전사이므로, $\sigma_i|_{U_{i, j}}$가 어떤 원소
$\tau_{i, j} \in \mathcal{G}(U_{i, j})$의 상이 되게 하는 덮개들
$\{U_{i, j} \to U_i\}$가 존재한다. 범주
$\text{SR}(\mathcal{C}, X)$의 대상 $Z = \{U_{i, j} \to X\}$와 그 명백한 사상
$u : Z \to K_p$를 생각하자. Definition
\ref{hypercovering-definition-covering-SR}에 의해 $u$는 분명히 덮개이다. Lemma
\ref{hypercovering-lemma-covering}에 의해 $L_p \to K_p$가 $u$를 거쳐 인수분해되는
$X$의 초덮개 사상 $L \to K$가 존재한다. 따라서 $K$를 $L$로 바꾼 뒤에는
$\sigma$가 어떤 원소 $\tau \in \mathcal{G}(K_p)$의 상이라고 가정해도 된다.
$d(\sigma) = 0$이지만 반드시 $d(\tau) = 0$인 것은 아님에 유의하자.
따라서 $d(\tau) \in \mathcal{F}(K_{p + 1})$는 코사이클이다.
이 상황에서 $\delta(\xi)$를 코사이클 $d(\tau)$의
$\check{H}^{p + 1}_{HC}(X, \mathcal{F})$에서의 류로 정의한다.

\medskip\noindent
이제 몇 가지를 확인해야 한다. (a) $\delta(\xi)$가 $\tau$의 선택에
의존하지 않는다는 것, (b) $\delta(\xi)$가 $\sigma$를 들어 올릴 수 있게 하는
초덮개 사상 $L \to K$의 선택에 의존하지 않는다는 것, 그리고
(c) $\delta(\xi)$가 $\xi$를 나타내기 위해 처음 택한 초덮개와
$\sigma$에 의존하지 않는다는 것이다. (a), (b), (c)의 확인은 생략한다.
초덮개의 선택과 무관하다는 것은 실제로 Lemmas
\ref{hypercovering-lemma-product-hypercoverings}와 \ref{hypercovering-lemma-homotopy}로 귀착된다.
또한 $\mathcal{C}$ 위의 아벨 층들의 짧은 완전열 사이의 사상에 관해
$\delta$가 함자적이라는 확인도 생략한다.

\medskip\noindent
끝으로 이 $\delta$의 정의를 사용하면 위의 아벨 층들의 짧은 완전열이
{\v C}ech 코호몰로지 군들의 긴 완전열을 낳는다는 것을 확인해야 한다.
먼저 $\delta(\xi) = 0$이면(여기서 $\xi$는 위와 같다) $\xi$가 어떤 원소
$\xi' \in \check{H}^p_{HC}(X, \mathcal{G})$의 상임을 보이자.
실제로 $\delta(\xi) = 0$이면, 위의 표기 아래 $d(\tau)$의 류는
$\check{H}^{p + 1}_{HC}(X, \mathcal{F})$에서 0이다. 따라서
$d(\tau)$를 $\mathcal{F}(L_{p + 1})$의 원소로 제한한 것이 어떤
$\upsilon \in \mathcal{F}(L_p)$에 대한 $d(\upsilon)$과 같아지게 하는
초덮개의 사상 $L \to K$가 존재한다. 이는
$\tau|_{L_p} + \upsilon$이 코사이클을 이루고, 원하는 대로 $\xi$로 가는
류 $\xi' \in \check{H}^p(L, \mathcal{G})$를 정한다는 뜻이다.

\medskip\noindent
$\xi' \in \check{H}^{p + 1}_{HC}(X, \mathcal{F})$가
$\check{H}^{p + 1}_{HC}(X, \mathcal{G})$에서 0으로 가면, 어떤
$\xi \in \check{H}^p_{HC}(X, \mathcal{H})$에 대해
그 원소가 $\delta(\xi)$와 같다는 것의 증명은 생략한다.
\end{proof}

\noindent
다음으로 약간의 기교를 써서 Theorem
\ref{hypercovering-theorem-cohomology-hypercoverings}의 Verdier 경우를 도출한다.

\begin{proposition}
\label{hypercovering-proposition-cohomology-hypercoverings}
$\mathcal{C}$를 섬유곱과 두 대상의 곱을 갖는 사이트라 하자.
$\mathcal{F}$를 $\mathcal{C}$ 위의 아벨 층이라 하고
$i \geq 0$이라 하자. 그러면 다음이 성립한다.
\begin{enumerate}
\item 모든 $\xi \in H^i(\mathcal{F})$에 대해, $\xi$가 표준 사상
$\check{H}^i(K, \mathcal{F}) \to H^i(\mathcal{F})$의 상에 속하게 하는
초덮개 $K$가 존재한다.
\item $K, L$이 초덮개이고 $\xi_K \in \check{H}^i(K, \mathcal{F})$,
$\xi_L \in \check{H}^i(L, \mathcal{F})$가 $H^i(\mathcal{F})$의 같은
원소로 가면, 초덮개 $M$과 사상들 $M \to K$, $M \to L$이 존재하여
$\xi_K$와 $\xi_L$이 $\check{H}^i(M, \mathcal{F})$의 같은 원소로 간다.
\end{enumerate}
다시 말해 집합론적 문제를 제외하면, $\mathcal{C}$ 위에서
$\mathcal{F}$의 코호몰로지 군들은 모든 초덮개에 걸친
$\mathcal{F}$의 {\v C}ech 코호몰로지 군들의 여극한이다.
\end{proposition}

\begin{proof}
이 결과는 Theorem \ref{hypercovering-theorem-cohomology-hypercoverings}의 자명한
귀결이다. 실제로 $\mathcal{C}$를 인위적으로 조금 더 큰 사이트
$\mathcal{C}'$로 바꿀 수 있는데, 이때 (I) $\mathcal{C}'$는 끝 대상
$X$를 갖고, (II) $\mathcal{C}$의 초덮개는 대체로
$\mathcal{C}'$에서 $X$의 초덮개와 같은 것이다. 그러나 사정의 성질상
상당한 정리 작업이 필요하다.

\medskip\noindent
$\mathcal{C}$에서 공통 끝점을 갖는 사상족 $\{U_i \to U\}$를, 사상
$\coprod_{i \in I} h_{U_i} \to h_U$가 층화 뒤 전사가 될 때
{\it 약한 덮개}라 하자. 새 사이트 $\mathcal{C}'$를 다음과 같이 구성한다.
\begin{enumerate}
\item 범주로서 $\Ob(\mathcal{C}') = \Ob(\mathcal{C}) \amalg \{X\}$로
두고, $\mathcal{C}'$의 각 대상에서 $X$로 가는 유일한 사상을 덧붙인다.
\item $\mathcal{C}$에 섬유곱과 두 대상의 곱이 존재하므로
$\mathcal{C}'$는 섬유곱을 갖는다.
\item $\mathcal{C}'$의 덮개는 $\mathcal{C}$의 약한 덮개들과 다음 조건을
만족하는 $\{U_i \to X\}_{i \in I}$들이다. 어떤 $i$에 대해
$U_i = X$이거나, 모든 $i$에 대해 $U_i \not = X$이고 $\mathcal{C}$ 위의
전층 사상 $\coprod h_{U_i} \to *$가 $\mathcal{C}$ 위에서 층화한 뒤
전사가 된다.
\item Sets, Lemma \ref{sets-lemma-coverings-site}를 적용하여 덮개들을
제한하고 사이트 $\mathcal{C}'$를 얻는다.
\end{enumerate}
포함 함자 $\mathcal{C} \to \mathcal{C}'$는 특수 쌍대연속 함자이므로
$\Sh(\mathcal{C}') = \Sh(\mathcal{C})$이다(Sites, Definition
\ref{sites-definition-special-cocontinuous-functor}를 보라).
직접적인 확인은 생략한다.

\medskip\noindent
모든 $i$에 대해 $U_i$가 $\mathcal{C}$의 대상이 되는
$\mathcal{C}'$의 덮개 $\{U_i \to X\}$를 택하자
($\mathcal{C} \to \mathcal{C}'$가 특수 쌍대연속이므로 가능하다).
그러면 $K_0 = \{U_i \to X\}$는 위에서 구성한 사이트
$\mathcal{C}'$의 덮개이다. $K_0$를 $\text{SR}(\mathcal{C}', X)$의
대상으로 보고 $K_{init} = \text{cosk}_0(K_0)$로 둔다. Example
\ref{hypercovering-example-cech}에 의해 $K_{init}$는 $X$의 초덮개이다. 모든
$K_{init, n}$은 $W_j \in \Ob(\mathcal{C})$인
$\{W_j \to X\}$ 꼴임에 유의하자.

\medskip\noindent
(1)의 증명. $\xi \in H^i(\mathcal{F}) = H^i(X, \mathcal{F}')$를 택하자.
여기서 $\mathcal{F}'$는 $\mathcal{C}$ 위의 $\mathcal{F}$에 대응하는
$\mathcal{C}'$ 위의 아벨 층이다. Theorem
\ref{hypercovering-theorem-cohomology-hypercoverings}에 의해, $\xi$가
$\check{H}^i(K', \mathcal{F})$의 한 원소에서 오게 하는
$\mathcal{C}'$에서 $X$의 초덮개 사상 $K' \to K_{init}$가 존재한다.
$K'_n = \{U_{n, j} \to X\}$라 쓰자. 이제 $K'_n$이 $K_{init, n}$으로
가므로 $U_{n, j}$는 $\mathcal{C}$의 대상이다. 따라서
$K_n = \{U_{n, j}\}$로 두어 $\text{SR}(\mathcal{C})$의 단체적 대상
$K$를 정의할 수 있다. $\mathcal{C}$의 사상족들로 이루어진
$\mathcal{C}'$의 덮개는 약한 덮개이므로, $K$는 Definition
\ref{hypercovering-definition-hypercovering-variant}의 의미에서 초덮개이다.
끝으로 $\mathcal{F}'$는 $\mathcal{C}'$ 위에서 그 $\mathcal{C}$로의
제한이 $\mathcal{F}$와 같은 유일한 층이므로, {\v C}ech 복합체들
$s(\mathcal{F}(K))$와 $s(\mathcal{F}'(K'))$는 동일하다. 따라서 (1)이
따른다. (코호몰로지 군으로 가는 사상과의 양립성은 생략한다.)

\medskip\noindent
(2)의 증명. $K$와 $L$을 $\mathcal{C}$의 초덮개라 하자.
$K'$와 $L'$을 함자
$\text{SR}(\mathcal{C}) \to \text{SR}(\mathcal{C}', X)$,
$\{U_i\} \mapsto \{U_i \to X\}$로 $K$와 $L$에서 얻는
$\text{SR}(\mathcal{C}', X)$의 단체적 대상들이라 하자. 위와 같이
{\v C}ech 복합체들이 같으므로, $\mathcal{C}'$ 위에서
$\mathcal{F}'$의 같은 코호몰로지 류로 가는 $\xi_{K'}$와
$\xi_{L'}$를 얻는다. 집합론적 문제 때문에 필요하다면
$\mathcal{C}'$에서 선택한 덮개들을 늘린 뒤, $K'$와 $L'$이
$\mathcal{C}'$에서 $X$의 초덮개라고 가정해도 된다. 이는 Definition
\ref{hypercovering-definition-hypercovering-variant}의 초덮개 정의와
$\mathcal{C}$의 약한 덮개가 $\mathcal{C}'$의 덮개를 준다는 사실에
의해 성립한다. Theorem \ref{hypercovering-theorem-cohomology-hypercoverings}에 의해
$\mathcal{C}'$에서 $X$의 초덮개 $M'$과 사상들
$M' \to K'$, $M' \to L'$, $M' \to K_{init}$가 존재하여,
$\xi_{K'}$와 $\xi_{L'}$가 $\check{H}^i(M', \mathcal{F})$의 같은
원소로 제한된다. 이 명제를 위와 같이 풀어 쓰면 (2)가 성립함을 얻는다.
\end{proof}





\section{공간의 초덮개}
\label{hypercovering-section-hypercoverings-spaces}

\noindent
위의 이론은 위상공간의 경우에도 어느 정도 흥미롭다. 이 경우에는
초덮개가 무엇인지 구체적으로 풀어 쓰고 그 결과가 실제로 무엇을
말하는지 볼 수 있다.

\medskip\noindent
$X$를 위상공간이라 하자. Sites, Example
\ref{sites-example-site-topological}의 사이트 $X_{Zar}$를 생각하자.
$X_{Zar}$의 대상은 단순히 $X$의 열린집합이고, $X_{Zar}$의 사상은
단순히 포함에 대응함을 상기하자. 그렇다면 사이트 $X_{Zar}$에서
$X$의 초덮개란 무엇인가?

\medskip\noindent
먼저 Definition \ref{hypercovering-definition-SR}을 풀어 쓰자.
$\text{SR}(X_{Zar}, X)$의 대상은 단순히 집합 $I$와 각
$i \in I$에 대한 열린집합 $U_i \subset X$로 주어진다.
사상 $U_i \to X$에 혼동의 여지가 없으므로 이를
$\{U_i\}_{i \in I}$로 쓰자. 이러한 두 대상 사이의 사상
$\{U_i\}_{i \in I} \to \{V_j\}_{j \in J}$는 모든 $i \in I$에 대해
$U_i \subset V_{\alpha(i)}$를 만족하는 집합의 사상
$\alpha : I \to J$로 주어진다. 이러한 사상이 덮개인 것은 언제인가?
모든 $j \in J$에 대해
$V_j = \bigcup_{i\in I, \ \alpha(i) = j} U_i$일 때 그리고 그때에만
그러하다(또한 이는 사이트 $X_{Zar}$의 덮개이다).

\medskip\noindent
위의 내용을 사용하면 사이트 $X_{Zar}$의 초덮개를 다음과 같이
기술할 수 있다. $X_{Zar}$에서 $X$의 초덮개는 다음 자료로 주어진다.
\begin{enumerate}
\item 단체 집합 $I$(Simplicial, Section
\ref{simplicial-section-simplicial-set}를 보라),
\item 각 $n \geq 0$과 모든 $i \in I_n$에 대한 열린집합 $U_i \subset X$.
\end{enumerate}
이러한 자료의 모음을 $(I, \{U_i\})$로 나타내자. 이것이 $X$의
초덮개가 되려면 다음 성질들을 요구한다.
\begin{itemize}
\item $i \in I_n$이고 $0 \leq a \leq n$이면
$U_i \subset U_{d^n_a(i)}$이다.
\item $i \in I_n$이고 $0 \leq a \leq n$이면 $U_i = U_{s^n_a(i)}$이다.
\item 다음이 성립한다.
\begin{equation}
\label{hypercovering-equation-covering-X}
X = \bigcup\nolimits_{i \in I_0} U_i,
\end{equation}
\item 모든 $i_0, i_1 \in I_0$에 대해 다음이 성립한다.
\begin{equation}
\label{hypercovering-equation-covering-two}
U_{i_0} \cap U_{i_1} =
\bigcup\nolimits_{i \in I_1, \ d^1_0(i) = i_0, \ d^1_1(i) = i_1} U_i,
\end{equation}
\item 모든 $n \geq 1$과, 모든 $0\leq a < b\leq n + 1$에 대해
$d^n_{b - 1}(i_a) = d^n_a(i_b)$를 만족하는 모든
$(i_0, \ldots, i_{n + 1}) \in (I_n)^{n + 2}$에 대해 다음이 성립한다.
\begin{equation}
\label{hypercovering-equation-covering-general}
U_{i_0} \cap \ldots \cap U_{i_{n + 1}} =
\bigcup\nolimits_{i \in I_{n + 1},
\ d^{n + 1}_a(i) = i_a, \ a = 0, \ldots, n + 1} U_i,
\end{equation}
\item 열린 덮개들 (\ref{hypercovering-equation-covering-X}),
(\ref{hypercovering-equation-covering-two}), (\ref{hypercovering-equation-covering-general})은 각각
$\text{Cov}(X_{Zar})$의 원소이다(이는 덮개의 첨자집합 크기에 상계를
주는 집합론적 조건이다).
\end{itemize}
예를 들어 조건들 (\ref{hypercovering-equation-covering-X})와
(\ref{hypercovering-equation-covering-two})는 공간 위의 층을 다룬 장에서 익숙할 것이고,
조건 (\ref{hypercovering-equation-covering-general})은 그 자연스러운 일반화이다.

\begin{remark}
\label{hypercovering-remark-not-covering-set}
이 기술의 한 특징은 다중 교집합들
$U_{i_0} \cap \ldots \cap U_{i_{n + 1}}$ 가운데 하나가 공집합이면
오른쪽의 덮개가 빈 덮개일 수 있다는 점이다. 따라서 사상들
$I_{n + 1} \to (\text{cosk}_n\text{sk}_n I)_{n + 1}$이 자동으로
전사가 되는 것은 아니다. 이는 $I$의 기하학적 실현이 흥미로운
(축약 가능하지 않은) 공간일 수 있음을 뜻한다.

\medskip\noindent
실제로 $I'_n \subset I_n$을 $U_i \not = \emptyset$을 만족하는 단체
$i \in I_n$들로 이루어진 부분집합이라 하자. $I' \subset I$가 단체적
부분집합이고 $(I', \{U_i\})$가 초덮개임은 쉽게 알 수 있다. 따라서
어떤 초덮개도 모든 열린집합 $U_i$가 공집합이 아닌 초덮개로 세분할 수 있다.
\end{remark}

\begin{remark}
\label{hypercovering-remark-repackage-into-simplicial-space}
이 정보를 다시 다른 방식으로 묶어 보자. $(I, \{U_i\})$가 위상공간
$X$의 초덮개라고 하자. 이 자료가 주어지면
$$
U_n = \coprod\nolimits_{i \in I_n} U_i,
$$
로 두어 단체적 위상공간 $U_\bullet$을 구성할 수 있다. 또한 주어진
$\varphi : [n] \to [m]$에 대해 사상 $U(\varphi) : U_n \to U_m$을
각 $i \in I_n$에 대한 포함 $U_i \subset U_{\varphi(i)}$에서 오는
사상으로 둔다. 이 단체적 위상공간에는 증대
$\epsilon : U_\bullet \to X$가 딸려 있다. 이 사상에 의해 단체적 공간
$U_\bullet$은 $X$의 초덮개가 되고, 이를 따라 \cite[Expos\'e Vbis]{SGA4}의
의미에서 코호몰로지 강하가 성립한다. 다시 말해
$H^n(U_\bullet, \epsilon^*\mathcal{F}) = H^n(X, \mathcal{F})$이다.
(단체적 공간 위의 코호몰로지와 그 용어로 정식화된 코호몰로지 강하에
관한 참고문헌을 나중에 여기에 삽입할 것.)
$\mathcal{F}$를 $X$ 위의 아벨 층이라 하자. 이 경우 Lemma
\ref{hypercovering-lemma-cech-spectral-sequence}의 스펙트럼 열은 다음 $E_1$-항을
갖는 스펙트럼 열이 된다.
$$
E_1^{p, q} = H^q(U_p, \epsilon_q^*\mathcal{F})
\Rightarrow
H^{p + q}(U_\bullet, \epsilon^*\mathcal{F}) = H^{p + q}(X, \mathcal{F})
$$
이는 $\epsilon^*\mathcal{F}$의 전체 코호몰로지를 $U_\bullet$의 각 조각
위에서 $\mathcal{F}$의 코호몰로지 군들과 비교한다. (이 스펙트럼 열에
관한 참고문헌을 나중에 여기에 삽입할 것.)
\end{remark}

\noindent
위상수학에서는 모든 $U_i$가 $X$의 위상의 주어진 기저에서 나오고,
모든 덮개 (\ref{hypercovering-equation-covering-two})와
(\ref{hypercovering-equation-covering-general})이 덮개들의 주어진 공종 모음에서
나온다는 성질을 갖는 $X$의 초덮개를 찾고 싶을 때가 많다.
다음은 두 가지 예시 보조정리이다.

\begin{lemma}
\label{hypercovering-lemma-basis-hypercovering}
$X$를 위상공간이라 하자. $\mathcal{B}$를 $X$의 위상의 기저라 하자.
각 $U_i$가 $\mathcal{B}$의 원소가 되는 $X$의 초덮개
$(I, \{U_i\})$가 존재한다.
\end{lemma}

\begin{proof}
$n \geq 0$이라 하자. $X$의 {\it $n$-절단 초덮개}란 $n$-절단
단체 집합 $I$와, 각 $i \in I_a$, $0 \leq a \leq n$에 대한 $X$의
열린집합 $U_i$로 주어지며, 초덮개를 정의하는 조건들이 의미가 있을
때마다 성립하는 것이라 하자. 다시 말해 그 안에 나타나는 모든 단체가
$a \leq n$인 $a$-단체일 때에만 포함 관계와 덮개 조건을 요구한다.
모든 $U_i \in \mathcal{B}$인 $n$-절단 초덮개 $(I, \{U_i\})$가
주어졌을 때, $a \leq n$인 $a$-단체를 하나도 더하지 않고 이를
$(n + 1)$-절단 초덮개로 확장할 수 있음을 보이면 보조정리가 따른다.
이는 다음과 같이 한다. 먼저 다음으로 정의되는 $(n + 1)$-절단
단체 집합 $I'$을 생각하자.
$I' = \text{sk}_{n + 1}(\text{cosk}_n I)$.
다음을 상기하자.
$$
I'_{n + 1} =
\left\{
\begin{matrix}
(i_0, \ldots, i_{n + 1}) \in (I_n)^{n + 2} \text{ such that}\\
d^n_{b - 1}(i_a) = d^n_a(i_b) \text{ for all }0\leq a < b\leq n + 1
\end{matrix}
\right\}
$$
$i' \in I'_{n + 1}$이 퇴화되어 있고 이를테면
$i' = s^n_a(i)$이면 $U_{i'} = U_i$로 둔다(어차피 두 번째 조건이
이를 강제한다). 이 경우 $J_{i'} = \{i'\}$로도 둔다.
$i' \in I'_{n + 1}$이 비퇴화이고 이를테면
$i' = (i_0, \ldots, i_{n + 1})$이면 집합 $J_{i'}$와 열린 덮개
\begin{equation}
\label{hypercovering-equation-choose-covering}
U_{i_0} \cap \ldots \cap U_{i_{n + 1}} =
\bigcup\nolimits_{i \in J_{i'}} U_i,
\end{equation}
를 택하되, 각 $i \in J_{i'}$에 대해 $U_i \in \mathcal{B}$가 되게 한다.
다음과 같이 두자.
$$
I_{n + 1} = \coprod\nolimits_{i' \in I'_{n + 1}} J_{i'}
$$
표준 사상 $\pi : I_{n + 1} \to I'_{n + 1}$이 있으며, 구성에 의해
$I'_{n + 1}$의 퇴화 단체들의 집합 위에서 전단사이다.
$i \in I_{n + 1}$에 대해 $d^{n + 1}_a(i) = d^{n + 1}_a(\pi(i))$로 정의한다.
$i \in I_n$에 대해 $s^n_a(i) \in I_{n + 1}$을 퇴화 단체
$s^n_a(i) \in I'_{n + 1}$ 위에 놓인 유일한 단체로 정의한다.
이것이 $X$의 $(n + 1)$-절단 초덮개를 정의한다는 확인은 생략한다.
\end{proof}

\begin{lemma}
\label{hypercovering-lemma-quasi-separated-quasi-compact-hypercovering}
$X$를 위상공간이라 하자. $\mathcal{B}$를 $X$의 위상의 기저라 하자.
다음을 가정하자.
\begin{enumerate}
\item $X$는 준콤팩트이다.
\item 각 $U \in \mathcal{B}$는 준콤팩트 열린집합이다.
\item $X$의 임의의 두 준콤팩트 열린집합의 교집합은 준콤팩트이다.
\end{enumerate}
그러면 다음 성질들을 갖는 $X$의 초덮개 $(I, \{U_i\})$가 존재한다.
\begin{enumerate}
\item 각 $U_i$는 기저 $\mathcal{B}$의 원소이다.
\item 각 $I_n$은 유한집합이고, 특히
\item 덮개들 (\ref{hypercovering-equation-covering-X}),
(\ref{hypercovering-equation-covering-two}), (\ref{hypercovering-equation-covering-general})은
각각 유한하다.
\end{enumerate}
\end{lemma}

\begin{proof}
Lemma \ref{hypercovering-lemma-basis-hypercovering}의 증명에 나온 구성에서
(\ref{hypercovering-equation-choose-covering})의 덮개를 $\mathcal{B}$의 원소들로
이루어진 유한 덮개로 택하면 곧바로 따른다. 세부 사항은 생략한다.
\end{proof}







\section{초덮개의 구성}
\label{hypercovering-section-hypercovering-sites}

\noindent
$\mathcal{C}$를 사이트라 하자. 이 절에서는
$\text{SR}(\mathcal{C})$의 단체적 대상을 다음과 같이 생각한다.
평소처럼 $K_n = K([n])$으로 놓고, $\varphi : [m] \to [n]$에 결부된
사상을 $K(\varphi) : K_n \to K_m$으로 나타낸다.
$K_n = \{U_{n, i}\}_{i \in I_n}$으로 쓸 수 있다.
$\varphi : [m] \to [n]$에 대해 사상 $K(\varphi) : K_n \to K_m$은
사상 $\alpha(\varphi) : I_n \to I_m$과, $i \in I_n$마다 주어진 사상
$f_{\varphi, i} : U_{n, i} \to U_{m, \alpha(\varphi)(i)}$로 정해진다.
$K$가 $\text{SR}(\mathcal{C})$의 단체적 대상이라는 사실로부터
$(I_n, \alpha(\varphi))$이 단체 집합이고, $\psi : [l] \to [m]$일 때
$f_{\psi, \alpha(\varphi)(i)} \circ f_{\varphi, i} =
f_{\varphi \circ \psi, i}$임을 알 수 있다.

\begin{lemma}
\label{hypercovering-lemma-split}
$\mathcal{C}$를 사이트라 하고, $K$를 $\text{SR}(\mathcal{C})$의
$r$-절단 단체적 대상이라 하자. 다음 조건들은 서로 동치이다.
\begin{enumerate}
\item $K$는 분할되어 있다(Simplicial, Definition
\ref{simplicial-definition-split}).
\item $f_{\varphi, i} : U_{n, i} \to U_{m, \alpha(\varphi)(i)}$
는 $r \geq n \geq 0$, 전사 $\varphi : [m] \to [n]$, $i \in I_n$에
대해 동형사상이다.
\item $f_{\sigma^n_j, i} : U_{n, i} \to U_{n + 1, \alpha(\sigma^n_j)(i)}$
는 $0 \leq j \leq n < r$, $i \in I_n$에 대해 동형사상이다.
\end{enumerate}
단체적 대상의 경우에도 같은 명제가 성립하며, 이때에는 (2)와 (3)에서
$r = \infty$로 놓으면 된다.
\end{lemma}

\begin{proof}
단체 집합의 분할은 유일하며 각 차수 $n$에서 비퇴화 첨자
$N(I_n)$으로 주어진다. Simplicial, Lemma
\ref{simplicial-lemma-splitting-simplicial-sets}를 보라.
$\text{SR}(\mathcal{C})$의 두 대상 $\{U_i\}_{i \in I}$와
$\{U_j\}_{j \in J}$의 쌍대곱은 명백한 표기법 아래
$\{U_l\}_{l \in I \amalg J}$로 주어진다. 따라서 $K$의 분할은
$N(K_n) = \{U_i\}_{i \in N(I_n)}$으로 주어져야 한다.
이제 정의를 풀어 쓰면 (1)과 (2)의 동치가 따른다. (2)와 (3)의 동치는
임의의 전사 $\varphi : [m] \to [n]$이
$k = n, n + 1, \ldots, m - 1$에 대한 사상 $\sigma^k_j$들의
합성이라는 사실에서 따른다.
\end{proof}

\begin{lemma}
\label{hypercovering-lemma-hypercovering-object}
$\mathcal{C}$를 섬유곱을 가지는 사이트라 하고,
$\mathcal{B} \subset \Ob(\mathcal{C})$를 부분집합이라 하자. 다음을 가정하자.
\begin{enumerate}
\item $\mathcal{C}$의 임의의 대상 $U$에는 $U_j \in \mathcal{B}$인
덮개 $\{U_j \to U\}_{j \in J}$가 있다.
\item $\{U_j \to U\}_{j \in J}$가 $U_j \in \mathcal{B}$인 덮개이고
$\{U' \to U\}$가 $U' \in \mathcal{B}$인 사상이면,
$\{U_j \to U\}_{j \in J} \amalg \{U' \to U\}$는 덮개이다.
\end{enumerate}
그러면 $\mathcal{C}$의 임의의 $X$에 대해 $X$의 초덮개 $K$가 존재하여,
$K_n = \{U_{n, i}\}_{i \in I_n}$이고 모든 $i \in I_n$에 대해
$U_{n, i} \in \mathcal{B}$이다.
\end{lemma}

\begin{proof}
Lemma \ref{hypercovering-lemma-basis-hypercovering}의 증명은 이 증명의 준비 연습이므로,
독자에게 먼저 그 증명을 읽기를 권한다.

\medskip\noindent
먼저 $\mathcal{C}$를 사이트 $\mathcal{C}/X$로 바꾼다.
그러면 $X$가 $\mathcal{C}$의 끝 대상이고 $\mathcal{C}$가 모든 유한 극한을
가진다고 가정할 수 있다(Categories, Lemma
\ref{categories-lemma-finite-limits-exist}).

\medskip\noindent
$n \geq 0$이라 하자. {\it $X$의 $n$-절단 $\mathcal{B}$-초덮개}란
$\text{SR}(\mathcal{C})$의 $n$-절단 단체적 대상 $K$로서,
$i \in I_a$, $0 \leq a \leq n$에 대해 $U_{a, i} \in \mathcal{B}$이고,
$K_0$가 $X$의 덮개이며,
$K_{a + 1} \to (\text{cosk}_a \text{sk}_a K)_{a + 1}$
이 $a = 0, \ldots, n - 1$에 대해 Definition
\ref{hypercovering-definition-covering-SR}의 의미에서 덮개인 것을 말한다.

\medskip\noindent
가정에 의해 $X$는 $U_i \in \mathcal{B}$인 덮개
$\{U_{0, i} \to X\}_{i \in I_0}$를 가지므로, $X$의 $0$-절단
$\mathcal{B}$-초덮개를 얻는다. $X$의 임의의 $0$-절단
$\mathcal{B}$-초덮개는 분할되어 있음에 유의하자. Lemma
\ref{hypercovering-lemma-split}를 보라.

\medskip\noindent
$n \geq 0$에 대해 $X$의 분할된 $n$-절단 $\mathcal{B}$-초덮개 $K$가
주어지면 이를 $X$의 분할된 $(n + 1)$-절단 $\mathcal{B}$-초덮개로
확장할 수 있음을 보이면 보조정리가 따른다.

\medskip\noindent
이 확장을 구성하자. $\text{SR}(\mathcal{C})$의 $(n + 1)$-절단
단체적 대상 $K' = \text{sk}_{n + 1}(\text{cosk}_n K)$를 생각하고 다음과 같이 쓰자.
$$
K'_{n + 1} = \{U'_{n + 1, i}\}_{i \in I'_{n + 1}}
$$
$K = \text{sk}_n K'$이므로 $0 \leq a \leq n$에 대해 $K_a = K'_a$이다.
모든 $i' \in I'_{n + 1}$에 대해 덮개
\begin{equation}
\label{hypercovering-equation-choose-covering-B}
\{g_{n + 1, j} : U_{n + 1, j} \to U'_{n + 1, i'}\}_{j \in J_{i'}}
\end{equation}
를 $j \in J_{i'}$에 대해 $U_{n + 1, j} \in \mathcal{B}$가 되도록 택한다.
이는 보조정리의 $\mathcal{B}$에 관한 가정으로 가능하다.
$0 \leq m \leq n$에 대해 비퇴화 첨자들의 부분집합을
$N_m \subset I_m$으로 나타내고 다음과 같이 놓는다.
$$
I_{n + 1} =
\coprod\nolimits_{\varphi : [n + 1] \to [m]\text{ surjective, }0\leq m \leq n}
N_m \amalg
\coprod\nolimits_{i' \in I'_{n + 1}} J_{i'}
$$
$j \in I_{n + 1}$에 대해 다음과 같이 놓는다.
$$
U_{n + 1, j} =
\left\{
\begin{matrix}
U_{m, i} & \text{이면} &
j = (\varphi, i) & \text{여기서} & \varphi : [n + 1] \to [m], i \in N_m \\
U_{n + 1, j} & \text{이면} &
j \in J_{i'} & \text{여기서} & i' \in I'_{n + 1}
\end{matrix}
\right.
$$
으로 놓으며 표기법은 명백하다.
$K_{n + 1} = \{U_{n + 1, j}\}_{j \in I_{n + 1}}$으로 놓자.
구성에 의해 모든 $j \in I_{n + 1}$에 대해 $U_{n + 1, j}$는
$\mathcal{B}$의 원소이다. 서로 양립하는 사상
$$
I_{n + 1} \to I'_{n + 1}
\quad\text{and}\quad
K_{n + 1} \to K'_{n + 1}
$$
을 정의하자. 구체적으로 첫 번째 사상은
$(\varphi, i) \mapsto \alpha'(\varphi)(i)$와
$(j \in J_{i'}) \mapsto i'$로 주어진다.
두 번째 사상에는 다음 사상들을 사용한다.
$$
f'_{\varphi, i} : U_{m, i} \to U'_{n + 1, \alpha'(\varphi)(i)}
\quad\text{and}\quad
g_{n + 1, j} : U_{n + 1, j} \to U'_{n + 1, i'}
$$
사상
$$
K_{n + 1} \to K'_{n + 1} =
(\text{cosk}_n \text{sk}_n K')_{n + 1} =
(\text{cosk}_n K)_{n + 1}
$$
은 Definition \ref{hypercovering-definition-covering-SR}의 의미에서 덮개라고 주장한다.
실제로 $i' \in I'_{n + 1}$이라 하자. $i'$이 비퇴화이면
$I_{n + 1}$에서 $i'$의 역상은 $J_{i'}$와 같고, 선택
(\ref{hypercovering-equation-choose-covering-B})에 의해 $U'_{n + 1, i'}$의 덮개를 얻는다.
반면 $i'$이 퇴화이면 유일한 쌍 $(\varphi, i)$에 대해
$I_{n + 1}$에서 $i'$의 역상은
$J_{i'} \amalg \{(\varphi, i)\}$이며, 선택
(\ref{hypercovering-equation-choose-covering-B})과 보조정리의 가정 (2)에 의해 덮개를 얻는다.

\medskip\noindent
증명을 마치려면 사상 $\varphi : [m] \to [n + 1]$,
$0 \leq m \leq n$에 대응하는 사상 $K(\varphi) : K_{n + 1} \to K_m$과,
사상 $\varphi : [n + 1] \to [m]$, $0 \leq m \leq n$에 대응하는 사상
$K(\varphi) : K_m \to K_{n + 1}$을 적절한 합성 관계를 만족하도록
정의해야 한다. 첫 번째 종류에는 합성
$$
K_{n + 1} \to K'_{n + 1} \xrightarrow{K'(\varphi)} K'_m = K_m
$$
을 사용하여 $K(\varphi) : K_{n + 1} \to K_m$을 정의한다.
두 번째 종류에 대해 $\varphi : [n + 1] \to [m]$,
$0 \leq m \leq n$이 주어졌다고 하자. 이에 대응하는 사상
$K(\varphi) : K_m \to K_{n + 1}$을 다음과 같이 정의한다.
\begin{enumerate}
\item $i \in I_m$에 대해 $\alpha(\psi)(i_0) = i$를 만족하는
유일한 전사 $\psi : [m] \to [m_0]$과 유일한 비퇴화 원소
$i_0 \in I_{m_0}$가 존재한다\footnote{예를 들어 $i$가 비퇴화이면
$m = m_0$이고 $\psi = \text{id}_{[m]}$이다.}.
\item $\varphi_0 = \psi_0 \circ \varphi : [n + 1] \to [m_0]$으로 놓고
$i \in I_m$을 $(\varphi_0, i_0) \in I_{n + 1}$로 보낸다. 다시 말해
$\alpha(\varphi)(i) = (\varphi_0, i_0)$으로 놓는다.
\item 사상
$f_{\varphi, i} : U_{m, i} \to U_{n + 1, \alpha(\varphi)(i)} = U_{m_0, i_0}$
은 동형사상 $f_{\psi, i_0} : U_{m_0, i_0} \to U_{m, i}$의 역이다
(Lemma \ref{hypercovering-lemma-split}를 보라).
\end{enumerate}
이것이 주어진 $n$-절단 초덮개를 확장하는 $X$의 분할된
$(n + 1)$-절단 $\mathcal{B}$-초덮개를 정의한다는 직접적이지만
번거로운 확인은 생략한다. 실제로 모든 $0 \leq a, b \leq n + 1$과
$\varphi : [a] \to [b]$에 대해 사상 $K(\varphi)$들이 올바르게
합성된다는 확인을 제외하면 위 구성에서 모든 것이 명백하다.
\end{proof}

\begin{lemma}
\label{hypercovering-lemma-hypercovering-site}
$\mathcal{C}$를 등화자와 섬유곱을 가지는 사이트라 하고,
$\mathcal{B} \subset \Ob(\mathcal{C})$를 부분집합이라 하자.
$\mathcal{C}$의 임의의 대상이 그 구성원들이 $\mathcal{B}$의 원소인
덮개를 가진다고 가정하자. 그러면 모든 $i \in I_n$에 대해
$U_i \in \mathcal{B}$인 $K_n = \{U_i\}_{i \in I_n}$을 만족하는
초덮개 $K$가 존재한다.
\end{lemma}

\begin{proof}
이 증명은 Lemma \ref{hypercovering-lemma-hypercovering-object}의 증명과 거의 같다.
차이점만 설명하겠다.

\medskip\noindent
$n \geq 1$이라 하자. {\it $n$-절단 $\mathcal{B}$-초덮개}란
$\text{SR}(\mathcal{C})$의 $n$-절단 단체적 대상 $K$로서,
$i \in I_a$, $0 \leq a \leq n$에 대해 $U_{a, i} \in \mathcal{B}$이고
다음을 만족하는 것을 말한다.
\begin{enumerate}
\item $F(K_0)^\# \to *$는 전사이다.
\item $F(K_1)^\# \to F(K_0)^\# \times F(K_0)^\#$는 전사이다.
\item $F(K_{a + 1})^\# \to F((\text{cosk}_a \text{sk}_a K)_{a + 1})^\#$
은 $a = 1, \ldots, n - 1$에 대해 전사이다.
\end{enumerate}
먼저 분할된 $1$-절단 $\mathcal{B}$-초덮개를 명시적으로 구성한다.

\medskip\noindent
$I_0 = \mathcal{B}$와 $K_0 = \{U\}_{U \in \mathcal{B}}$로 놓자.
$\mathcal{B}$에 관한 가정으로 (1)이 성립한다. 다음과 같이 놓는다.
$$
\Omega =
\{(U, V, W, a, b) \mid U, V, W \in \mathcal{B}, a : U \to V, b : U \to W\}
$$
이제 $I_1 = I_0 \amalg \Omega$로 놓는다. $i \in I_1$에 대해
$i \in I_0$이면 $U_{1, i} = U_{0, i}$로 놓고,
$i = (U, V, W, a, b) \in \Omega$이면 $U_{1, i} = U$로 놓는다.
사상 $K(\sigma^0_0) : K_0 \to K_1$은 포함
$\alpha(\sigma^0_0) : I_0 \to I_1$과 대상들 위의 항등사상
$f_{\sigma^0_0, i} : U_{0, i} \to U_{1, i}$에 대응한다.
사상 $K(\delta^1_0), K(\delta^1_1) : K_1 \to K_0$은
$I_0 \subset I_1$ 위에서 항등이고
$(U, V, W, a, b) \in \Omega \subset I_1$을 각각 $V$, $W$로 보내는
두 사상 $I_1 \to I_0$에 대응한다. 이에 대응하는 사상
$f_{\delta^1_0, i}, f_{\delta^1_1, i} : U_{1, i} \to U_{0, i}$는
$i \in I_0$이면 항등사상이고,
$i = (U, V, W, a, b) \in \Omega$이면 각각 $a,b$이다.
(2)가 성립하는 이유는 다음과 같다. $\mathcal{C}$의 대상 $U$ 위에서
$F(K_0)^\# \times F(K_0)^\#$의 임의의 단면은 $U$를 어떤 덮개의
구성원들로 바꾸고 나면 사상 $U \to F(K_0) \times F(K_0)$에서 온다.
이는 다시 $V, W \in \mathcal{B}$와 두 사상 $U \to V$, $U \to W$가
있음을 뜻한다. $U$를 다시 어떤 덮개의 구성원들로 바꾸면 원하는 대로
$U \in \mathcal{B}$라고 가정할 수 있다.

\medskip\noindent
$n \geq 1$에 대해 분할된 $n$-절단 $\mathcal{B}$-초덮개 $K$가
주어지면 이를 분할된 $(n + 1)$-절단 $\mathcal{B}$-초덮개로
확장할 수 있음을 보이면 보조정리가 따른다. 여기서 논증은
Lemma \ref{hypercovering-lemma-hypercovering-object}의 증명과 정확히 같은 방식으로
진행된다. 다음 설명을 제외한 정확한 세부 사항은 생략한다.
첫째, 사상 $K_{n + 1} \to (\text{cosk}_n K)_{n + 1}$ 자체가
덮개일 필요가 없으므로 현재 보조정리의 증명에서는 가정 (2)가
필요하지 않다. 이 사상이 상응하는 집합의 층들 위에서 전사를
유도하기만 하면 충분하며, 이는 Sites, Lemma
\ref{sites-lemma-covering-surjective-after-sheafification}에서 따른다.
둘째, $\mathcal{C}$가 섬유곱과 등화자를 가진다는 가정은
$\text{SR}(\mathcal{C})$도 섬유곱과 등화자를 가지며 $F$가 이들과
가환함을 보장한다(Lemma \ref{hypercovering-lemma-coprod-prod-SR}). 이는 사용한
여골격 함자들이 존재함을 보장하기에 충분하다(Simplicial, Remark
\ref{simplicial-remark-existence-cosk}와 Categories, Lemma
\ref{categories-lemma-fibre-products-equalizers-exist}를 보라).
\end{proof}

\begin{lemma}
\label{hypercovering-lemma-hypercovering-morphism-sites}
$f : \mathcal{C} \to \mathcal{D}$를 함자
$u : \mathcal{D} \to \mathcal{C}$로 주어지는 사이트의 사상이라 하자.
$\mathcal{D}$와 $\mathcal{C}$가 등화자와 섬유곱을 가지며 $u$가 이들과
가환한다고 가정하자. $\text{SR}(\mathcal{D})$의 단체적 대상 $K$가
초덮개이면 $u(K)$도 초덮개이다.
\end{lemma}

\begin{proof}
이 절의 도입부에서와 같이 $K_n = \{U_{n, i}\}_{i \in I_n}$으로 쓰면,
$u(K)$는 $u(K_n) = \{u(U_i)\}_{i \in I_n}$으로 주어지는
$\text{SR}(\mathcal{C})$의 대상이다. Sites, Lemma
\ref{sites-lemma-pullback-representable-sheaf}에 의해
$U \in \Ob(\mathcal{D})$에 대해 $f^{-1}h_U^\# = h_{u(U)}^\#$이다.
이는 모든 $n$에 대해 $f^{-1}F(K_n)^\# = F(u(K_n))^\#$임을 뜻한다.
$u(K)$가 초덮개가 되기 위한 Definition
\ref{hypercovering-definition-hypercovering-variant}의 조건 (1), (2), (3)을 확인하자.
$f^{-1}$가 완전 함자이므로
$$
F(u(K_0))^\# = f^{-1}F(K_0)^\# \to f^{-1}* = *
$$
는 전사인 사상의 당김이므로 전사이고, (1)을 얻는다. 마찬가지로
$$
F(u(K_1))^\# = f^{-1}F(K_1)^\# \to
f^{-1} (F(K_0) \times F(K_0))^\# = F(u(K_0))^\# \times F(u(K_0))^\#
$$
는 당김이므로 전사이고, (2)를 얻는다. 조건 (3)에 대해 같은 방법으로
결론을 얻으려면 다음이 성립하면 충분하다.
$$
F((\text{cosk}_n \text{sk}_n u(K))_{n + 1})^\# =
f^{-1}F((\text{cosk}_n \text{sk}_n K)_{n + 1})^\#
$$
위 논의는 $f^{-1}F(-) = F(u(-))$임을 보인다. 따라서 $n \geq 1$에 대해
$u$가 $(\text{cosk}_n \text{sk}_n K)_{n + 1}$을 정의할 때 쓰이는
극한들과 가환함을 보이면 충분하다. Simplicial, Remark
\ref{simplicial-remark-existence-cosk}에 의해 이 극한들은 유한 연결
극한들이며, 가정에 의해 $u$는 이들과 가환한다.
\end{proof}

\begin{lemma}
\label{hypercovering-lemma-hypercovering-continuous-functor}
$\mathcal{C}$와 $\mathcal{D}$를 사이트라 하고
$u : \mathcal{D} \to \mathcal{C}$를 연속 함자라 하자.
$\mathcal{D}$와 $\mathcal{C}$가 섬유곱을 가지며 $u$가 이들과
가환한다고 가정하자. $Y \in \mathcal{D}$라 하고
$K \in \text{SR}(\mathcal{D}, Y)$를 $Y$의 초덮개라 하자.
그러면 $u(K)$는 $u(Y)$의 초덮개이다.
\end{lemma}

\begin{proof}
대상의 초덮개라는 개념이 더 강하므로 이는 Lemma
\ref{hypercovering-lemma-hypercovering-morphism-sites}의 증명보다 쉽다. Definitions
\ref{hypercovering-definition-hypercovering}과 \ref{hypercovering-definition-covering-SR}을 보라.
실제로 사이트의 사상의 정의에 의해 $u$는 덮개를 덮개로 보낸다.
$n \geq 1$에 대해 $(\text{cosk}_n \text{sk}_n K)_{n + 1}$을 정의할 때
쓰이는 극한들과 $u$가 가환함을 확인하면 충분하다. 유도된 함자
$\mathcal{D}/Y \to \mathcal{C}/X$가 모든 유한 극한과 가환하므로 이는
명백하다(그리고 Categories, Lemma
\ref{categories-lemma-finite-limits-exist}에 의해 정의역과 공역은
모든 유한 극한을 가진다).
\end{proof}

\begin{lemma}
\label{hypercovering-lemma-w-contractible}
$\mathcal{C}$를 사이트라 하고, $\mathcal{B} \subset \Ob(\mathcal{C})$를
부분집합이라 하자. 다음을 가정하자.
\begin{enumerate}
\item $\mathcal{C}$는 섬유곱을 가진다.
\item 모든 $X \in \Ob(\mathcal{C})$에 대해 $U_i \in \mathcal{B}$인
유한 덮개 $\{U_i \to X\}_{i \in I}$가 존재한다.
\item $\{U_i \to X\}_{i \in I}$가 $U_i \in \mathcal{B}$인 유한 덮개이고
$U \to X$가 $U \in \mathcal{B}$인 사상이면,
$\{U_i \to X\}_{i \in I} \amalg \{U \to X\}$는 덮개이다.
\end{enumerate}
그러면 모든 $X$에 대해 $X$의 초덮개 $K$가 존재하여, 각
$K_n = \{U_{n, i} \to X\}_{i \in I_n}$에서 $I_n$은 유한하고
$U_{n, i} \in \mathcal{B}$이다.
\end{lemma}

\begin{proof}
이 보조정리는 사이트에 대한 Lemma
\ref{hypercovering-lemma-quasi-separated-quasi-compact-hypercovering}의 유사 명제이다.
다음 두 가지에 유의하면서 Lemma
\ref{hypercovering-lemma-hypercovering-object}의 증명을 그대로 따르면 된다.
\begin{enumerate}
\item[(a)] 처음 덮개 $\{U_{0, i} \to X\}_{i \in I_0}$를
$U_{0, i} \in \mathcal{B}$이고 첨자 집합 $I_0$가 유한하도록 택한다.
\item[(b)] 덮개 (\ref{hypercovering-equation-choose-covering-B})를 택할 때
$J_{i'}$를 유한하게 택한다.
\end{enumerate}
이렇게 바꾸면 모든 $n$에 대해 유한 첨자 집합 $I_n$을 얻게 됨은
쉽게 알 수 있다.
\end{proof}

\begin{remark}
\label{hypercovering-remark-taking-disjoint-unions}
$\mathcal{C}$를 사이트라 하고, $K$와 $L$을
$\text{SR}(\mathcal{C})$의 대상이라 하자.
$K = \{U_i\}_{i \in I}$와 $L = \{V_j\}_{j \in J}$로 쓰자.
$U = \coprod_{i \in I} U_i$와 $V = \coprod_{j \in J} V_j$가
존재한다고 가정하자. 그러면 다음을 얻는다.
$$
\Mor_{\text{SR}(\mathcal{C})}(K, L) \longrightarrow \Mor_\mathcal{C}(U, V)
$$
구성은 다음과 같다. $\alpha : I \to J$와
$f_i : U_i \to V_{\alpha(i)}$로 주어진 $f : K \to L$에 대해 함자들의 변환
$$
\Mor_\mathcal{C}(V, -) =
\prod\nolimits_{j \in J} \Mor_\mathcal{C}(V_j, -)
\to
\prod\nolimits_{i \in I} \Mor_\mathcal{C}(U_i, -) =
\Mor_\mathcal{C}(U, -)
$$
을 얻는데, 이는 $(g_j)_{j \in J}$를
$(g_{\alpha(i)} \circ f_i)_{i \in I}$로 보낸다. 따라서 요네다 보조정리는
이에 대응하는 사상 $U \to V$를 준다. 물론 $U \to V$는 사상 $f_i$를 통해
쌍대곱 성분 $U_i$를 쌍대곱 성분 $V_{\alpha(i)}$로 보낸다.
\end{remark}

\begin{remark}
\label{hypercovering-remark-take-unions-hypercovering}
$\mathcal{C}$를 사이트라 하자. $\mathcal{C}$가 섬유곱과 등화자를
가진다고 가정하고 $K$를 초덮개라 하자.
$K_n = \{U_{n, i}\}_{i \in I_n}$으로 쓰자. 다음을 가정하자.
\begin{enumerate}
\item[(a)] $U_n = \coprod_{i \in I_n} U_{n, i}$가 존재한다.
\item[(b)] $\coprod_{i \in I_n} h_{U_{n, i}} \to h_{U_n}$는
층화한 뒤 동형사상을 유도한다.
\end{enumerate}
그러면 $L_n = \{U_n\}$인 $\text{SR}(\mathcal{C})$의 또 다른 단체적 대상
$L$을 얻는다. Remark \ref{hypercovering-remark-taking-disjoint-unions}를 보라.
이제 $L$이 초덮개라고 주장한다. 이를 위해 Definition
\ref{hypercovering-definition-hypercovering-variant}의 조건 (1), (2), (3)을 확인한다.
조건 (1)은 (b)와 $K$에 대한 (1)에서 따르고, 조건 (2)도 정확히 같은
방식으로 따른다. 조건 (3)은 다음 등식에서 따른다.
\begin{align*}
F((\text{cosk}_n \text{sk}_n L)_{n + 1})^\#
& =
((\text{cosk}_n \text{sk}_n F(L)^\#)_{n + 1}) \\
& =
((\text{cosk}_n \text{sk}_n F(K)^\#)_{n + 1}) \\
& =
F((\text{cosk}_n \text{sk}_n K)_{n + 1})^\#
\end{align*}
여기서 $n \geq 1$이다. 따라서 (1)과 (2)에서와 정확히 같은 방식으로
$K$에 대한 조건이 $L$에 대한 조건을 함의한다.
$F$가 연결 극한과 가환하고 층화가 완전하다는 사실이 첫 번째와 마지막
등식을 보인다. 가운데 등식은 (b)에 의해 $F(K)^\# = F(L)^\#$인 데서 따른다.
\end{remark}

\begin{remark}
\label{hypercovering-remark-take-unions-hypercovering-X}
$\mathcal{C}$를 사이트라 하고 $X \in \Ob(\mathcal{C})$라 하자.
$\mathcal{C}$가 섬유곱을 가진다고 가정하고 $K$를 $X$의 초덮개라 하자.
$K_n = \{U_{n, i}\}_{i \in I_n}$으로 쓰자. 다음을 가정하자.
\begin{enumerate}
\item[(a)] $U_n = \coprod_{i \in I_n} U_{n, i}$가 존재한다.
\item[(b)] $\text{SR}(\mathcal{C})$의 사상
$(\alpha, f_i) : \{U_i\}_{i \in I} \to \{V_j\}_{j \in J}$과
$(\beta, g_k) : \{W_k\}_{k \in K} \to \{V_j\}_{j \in J}$
가 주어지고 $U = \coprod U_i$, $V = \coprod V_j$, $W = \coprod W_j$가
존재하면 $U \times_V W =
\coprod_{(i, j, k), \alpha(i) = j = \beta(k)} U_i \times_{V_j} W_k$이다.
\item[(c)] $(\alpha, f_i) : \{U_i\}_{i \in I} \to \{V_j\}_{j \in J}$
가 Definition \ref{hypercovering-definition-covering-SR}의 의미에서 덮개이고,
$U = \coprod U_i$와 $V = \coprod V_j$가 존재하면 Remark
\ref{hypercovering-remark-taking-disjoint-unions}의 이에 대응하는 사상 $U \to V$는
$\mathcal{C}$의 덮개이다.
\end{enumerate}
그러면 $L_n = \{U_n\}$인 $\text{SR}(\mathcal{C})$의 또 다른 단체적 대상
$L$을 얻는다. Remark \ref{hypercovering-remark-taking-disjoint-unions}를 보라.
이제 $L$이 $X$의 초덮개라고 주장한다. 이를 위해 Definition
\ref{hypercovering-definition-hypercovering}의 조건 (1), (2)를 확인한다.
조건 (1)은 (c)와 $K$에 대한 (1)에서 따른다. 실제로 $K$에 대한 (1)은
$K_0 = \{U_{0, i}\}_{i \in I_0}$가 Definition
\ref{hypercovering-definition-covering-SR}의 의미에서 $\{X\}$의 덮개라고 말한다.
조건 (2)가 따르는 이유는 $\mathcal{C}/X$가 모든 유한 극한을 가지므로
$\text{SR}(\mathcal{C}/X)$도 모든 유한 극한을 가지며, 조건 (b)가
``서로소 합을 취하는'' 구성이 이 유한 극한들과 가환한다고 말하기 때문이다.
따라서 사상
$$
L_{n + 1} \longrightarrow (\text{cosk}_n \text{sk}_n L)_{n + 1}
$$
은 ``서로소 합을 취하는'' 함자를 사상
$$
K_{n + 1} \longrightarrow (\text{cosk}_n \text{sk}_n K)_{n + 1}
$$
에 적용한 결과이므로 덮개이다. 뒤의 사상은 $K$에 대한 (2)에 의해
Definition \ref{hypercovering-definition-covering-SR}의 의미에서 덮개라고 가정되어 있다.
이는 의미가 있다. 실제로 성질 (b)는 특히 서로소 합을 취할 수 있는
$X$ 위의 반표현 가능 대상들로 이루어진 유한 도식에서 시작하면,
그 도식의 $\text{SR}(\mathcal{C}/X)$에서의 극한도 여전히 서로소 합을
취할 수 있는 $X$ 위의 반표현 가능 대상임을 보장한다.
\end{remark}
